% Matematica finanziaria applicata all'acquisto di un immobile residenziale.
% Documento di accompagnamento al progetto real-estate: formalizza ogni calcolo
% che il modello esegue, partendo dalle definizioni e senza passaggi impliciti.
% Compilazione: scripts/build.ps1 -Main docs\matematica-finanziaria.tex
\documentclass[11pt,a4paper]{article}

\usepackage[utf8]{inputenc}
\usepackage[T1]{fontenc}
\usepackage[italian]{babel}
\usepackage[a4paper,margin=2.6cm]{geometry}
\usepackage{amsmath}
\usepackage{amssymb}
\usepackage{amsthm}
\usepackage{booktabs}
\usepackage{enumitem}
\usepackage{float}
\usepackage{xcolor}
\usepackage{caption}
\usepackage[hidelinks]{hyperref}
\usepackage{listings}

\lstset{
  basicstyle=\ttfamily\small,
  breaklines=true,
  frame=single,
  framesep=4pt,
  rulecolor=\color{gray!50},
  showstringspaces=false,
  captionpos=b,
}

\theoremstyle{definition}
\newtheorem{definizione}{Definizione}[section]
\newtheorem{osservazione}[definizione]{Osservazione}
\theoremstyle{plain}
\newtheorem{proposizione}[definizione]{Proposizione}

\newcommand{\euro}{\texteuro}
\DeclareMathOperator{\VAN}{VAN}
\DeclareMathOperator{\TIR}{TIR}

\title{Matematica finanziaria dell'acquisto immobiliare\\[0.3em]
\large Formalizzazione completa dei calcoli del modello}
\author{Documento di accompagnamento al progetto \texttt{real-estate}}
\date{Revisione del 2 settembre 2026}

\begin{document}
\maketitle

\begin{abstract}
\noindent Questo documento formalizza la matematica che il modello di valutazione immobiliare esegue, dalla definizione di capitalizzazione fino alla graduatoria fra immobili, senza dare per scontato alcun passaggio. Ogni formula implementata nel codice compare qui con la sua derivazione, la sua notazione, le sue ipotesi e i suoi limiti; ogni scelta fra due forme algebricamente vicine ma numericamente diverse è motivata. L'ultimo capitolo contiene una tavola di corrispondenza fra i simboli usati qui, i nomi definiti nel workbook Excel e le funzioni Python del motore, così che ogni numero del modello sia rintracciabile fino alla sua formula. Il caso di riferimento del progetto è svolto per intero in appendice, con tutti i valori intermedi.

Il documento non è una dispensa di matematica finanziaria generale: è la matematica \emph{di questo problema}. Ciò che serve viene derivato, ciò che non serve non compare, e dove il dominio impone una complicazione — un'imposta a tratti, un'indicizzazione parziale, una convenzione contrattuale italiana — la complicazione è nel testo e non nascosta in una nota.
\end{abstract}

\tableofcontents
\newpage

\section{Come si legge questo documento, se non si e' abituati alle formule}

Questa sezione esiste perche' il resto del documento sarebbe illeggibile senza. Non spiega finanza: spiega i segni. Chi ha familiarita' con la notazione matematica puo' saltarla, e chi non ce l'ha dovrebbe leggerla per intero, perche' dopo di essa ogni formula del documento si legge come una frase.

\subsection{Una lettera e' il nome di un numero}

Quando si scrive $P = 120\,000$ si sta dando un nome a un numero, esattamente come si fa con una cella di un foglio di calcolo a cui si assegna un'etichetta. La lettera $P$ non e' un'incognita da trovare: e' un'abbreviazione. Si usa una lettera invece di scrivere ogni volta la parola prezzo per una ragione pratica: le formule diventano abbastanza corte da vedersi tutte insieme, e vedere tutta una formula insieme e' ciò che permette di accorgersi se ha senso.

La scelta delle lettere non e' casuale nel documento: la tabella della sezione seguente le elenca tutte, e conviene tenerla aperta accanto le prime volte. Le lettere greche indicano tipicamente delle quote o delle aliquote, cioe' numeri fra zero e uno: $\pi$ e' l'inflazione, $\theta$ un'aliquota d'imposta, $\rho$ una rivalutazione. Le lettere latine indicano importi in euro o conteggi: $P$ il prezzo, $C$ il costo totale, $n$ un numero di rate.

\subsection{Il pedice e' un'etichetta che dice quale}

Quando si scrive $D_k$ si intende il debito residuo dopo la rata numero $k$. Il numeretto in basso, il pedice, non moltiplica e non eleva: e' un'etichetta. $D_0$ e' il debito all'inizio, $D_1$ dopo la prima rata, $D_{300}$ dopo la trecentesima. Se in un foglio di calcolo il debito residuo occupa una colonna, allora $D_k$ e' semplicemente la cella di quella colonna alla riga $k$.

Un pedice puo' anche essere una parola abbreviata invece di un numero, e allora distingue due grandezze parenti: $L$ e' il canone potenziale e $L_e$ e' il canone effettivo, quello che entra davvero dopo lo sfitto. La $e$ non e' un numero, e' l'iniziale di effettivo.

\subsection{L'esponente e' una moltiplicazione ripetuta}

$(1+i)^n$ significa: prendi il numero $(1+i)$ e moltiplicalo per se stesso $n$ volte. Con $i = 0{,}03$ e $n = 2$ vale $1{,}03 \times 1{,}03 = 1{,}0609$.

Questo e' il segno piu' importante di tutto il documento, e vale capire perche' compare cosi' spesso. Un capitale che rende il tre per cento all'anno non cresce del tre per cento moltiplicato per gli anni, perche' nel secondo anno rende il tre per cento anche su quello che ha guadagnato nel primo. Dopo due anni non e' cresciuto del sei per cento ma del $6{,}09$, e dopo venticinque non del settantacinque per cento ma del centonove. La differenza fra moltiplicare e elevare e' tutta l'aritmetica dell'interesse composto, e in un orizzonte di venticinque anni non e' un dettaglio: e' il fenomeno principale.

L'esponente negativo, che compare come $(1+s)^{-n}$, significa dividere invece di moltiplicare: $(1+s)^{-n}$ e' un altro modo di scrivere $1/(1+s)^n$. Serve a fare il cammino inverso, cioe' a chiedersi quanto vale oggi un importo che arrivera' fra $n$ anni.

\subsection{Il segno di somma}

Il segno $\sum$, che e' la lettera greca sigma maiuscola, e' un'istruzione: somma tanti addendi costruiti tutti con la stessa regola. Si legge da tre punti.

\[
\underbrace{\sum_{k=1}^{4}}_{\text{per k che va da 1 a 4}} \underbrace{\frac{100}{(1{,}05)^k}}_{\text{somma questo}}
\;=\;
\frac{100}{1{,}05} + \frac{100}{1{,}05^2} + \frac{100}{1{,}05^3} + \frac{100}{1{,}05^4}
\;=\; 354{,}60
\]

Sotto il segno c'e' da dove parte il contatore, sopra dove arriva, a destra la cosa da sommare scritta una volta sola in funzione del contatore. E' esattamente cio' che fa un foglio di calcolo quando si scrive una formula in una cella, la si trascina in basso per quattro righe e poi si somma la colonna: la somma con il sigma e' quella stessa operazione, scritta in una riga invece che in cinque celle.

Nel documento il contatore si chiama quasi sempre $k$ quando conta rate o periodi e $t$ quando conta anni. Non c'e' nessuna differenza matematica fra le due lettere: e' una convenzione per far capire a colpo d'occhio di che periodo si stia parlando.

\subsection{La parentesi graffa che divide i casi}

Quando una grandezza si calcola in modi diversi a seconda della situazione, si scrive con una graffa e un elenco di casi:

\[
B(P) =
\begin{cases}
V_c, & \text{se si opta per il prezzo-valore},\\
P, & \text{altrimenti.}
\end{cases}
\]

Si legge cosi': la base imponibile $B$, che dipende dal prezzo $P$, vale $V_c$ nel primo caso e $P$ nel secondo. E' la stessa cosa che in un foglio di calcolo si scrive con un \texttt{SE}, e in un programma con un \texttt{if}. Le condizioni sono a destra, i valori a sinistra.

\subsection{$f(x)$ non e' una moltiplicazione}

Quando si scrive $C(P)$ non si sta moltiplicando $C$ per $P$: si sta dicendo che la grandezza $C$, il costo totale, dipende da $P$, il prezzo. La parentesi indica dipendenza, non prodotto. Serve quando la stessa grandezza va guardata a prezzi diversi, che e' esattamente il problema del capitolo sull'inversione: si scrive $C(P)$ perche' bisogna poter chiedere quanto vale il costo totale se il prezzo fosse un altro.

Quando invece si intende una moltiplicazione, il documento la scrive con il punto, come in $R \cdot \kappa$, oppure accostando i simboli, come in $P\,p$. L'accostamento e' la convenzione usuale in matematica e non e' ambigua perche' le parentesi sono riservate alla dipendenza.

\subsection{Massimo, minimo e valore assoluto}

$\max(a,b)$ e' il piu' grande fra i due numeri, $\min(a,b)$ il piu' piccolo. Compaiono perche' le norme fiscali sono spesso scritte in quella forma: l'imposta di registro e' la maggiore fra l'aliquota applicata alla base e un minimo di legge, cioe' $\max(B\,\tau,\,T_{\min})$, che si legge come prendi l'aliquota sulla base, ma se viene meno del minimo di legge paghi il minimo.

Il valore assoluto, scritto $|x|$, e' il numero senza il suo segno: $|-3| = 3$. Serve quando interessa quanto una cosa cambia e non in che direzione, per esempio nell'ordinare le variabili di un diagramma a tornado.

\subsection{Che cosa sono una proposizione e una dimostrazione}

Il documento contiene alcuni riquadri intitolati Proposizione, seguiti da un blocco che inizia con Dimostrazione. Non sono decorazioni accademiche e vale spiegare a cosa servono in un documento operativo.

Una proposizione e' un'affermazione precisa che si sostiene essere vera, per esempio che la rata di un mutuo alla francese vale una certa espressione. La dimostrazione e' il ragionamento che la sostiene, scritto in modo che chi legge possa verificarlo passo per passo invece di fidarsi. La differenza rispetto a una formula presentata senza dimostrazione e' pratica e non filosofica: una formula che si prende per buona si puo' solo confrontare con un'altra fonte, mentre una formula dimostrata si puo' controllare da soli, e se il ragionamento ha un buco lo si vede.

Il simbolo che chiude una dimostrazione, un quadratino, significa soltanto qui finisce il ragionamento.

\subsection{Un consiglio di lettura}

Nessuna formula di questo documento va imparata a memoria e nessuna va accettata sulla fiducia. Ognuna e' seguita da un paragrafo intitolato In parole che la rilegge come una frase italiana, di norma con un numero concreto preso dal caso di riferimento del progetto. Se una formula non torna, il paragrafo In parole dice cosa dovrebbe dire; se non torna nemmeno quello, la formula e' sbagliata o e' spiegata male, e in entrambi i casi vale segnalarlo.

\section{Notazione e convenzioni}

Si fissano una volta per tutte i simboli, perché la letteratura ne usa di incompatibili e l'ambiguità in questo dominio costa cifre.

\begin{table}[H]
\centering
\small
\begin{tabular}{@{}lll@{}}
\toprule
Simbolo & Significato & Unità \\
\midrule
$P$ & prezzo di acquisto trattato & \euro \\
$R$ & rendita catastale & \euro \\
$V_c$ & valore catastale, base del prezzo-valore & \euro \\
$T_a$ & imposte sul trasferimento & \euro \\
$A$ & costi accessori complessivi & \euro \\
$C$ & costo totale dell'operazione, $C = P + A$ & \euro \\
$M$ & capitale erogato dal mutuo & \euro \\
$E$ & capitale proprio immobilizzato, $E = C - M$ & \euro \\
$i$ & tasso annuo nominale del mutuo & adimensionale \\
$j$ & tasso di periodo, mensile & adimensionale \\
$n$ & numero di rate & conteggio \\
$N$ & orizzonte di analisi in anni & conteggio \\
$\Gamma$ & rata periodica del mutuo & \euro \\
$D_k$ & debito residuo dopo la rata $k$ & \euro \\
$L$ & canone annuo potenziale a pieno regime & \euro \\
$L_e$ & ricavo effettivo, al netto di sfitto e morosità & \euro \\
$\mathrm{NOI}$ & reddito operativo netto & \euro/anno \\
$U$ & utile netto annuo, dopo imposta sul canone & \euro/anno \\
$F_t$ & flusso di cassa dell'anno $t$ & \euro \\
$\pi$ & tasso di inflazione atteso & adimensionale \\
$g$ & tasso di indicizzazione del canone & adimensionale \\
$\rho$ & rivalutazione nominale annua dell'immobile & adimensionale \\
$s$ & tasso di sconto & adimensionale \\
$\theta$ & aliquota d'imposta sul canone & adimensionale \\
\bottomrule
\end{tabular}
\caption{Simboli usati nel documento.}
\end{table}

Tre convenzioni valgono in tutto il testo. I tassi sono espressi come frazioni, quindi il tre virgola due per cento è $0{,}032$ e non $3{,}2$; i flussi in uscita per chi compra sono negativi e quelli in entrata positivi, con l'unica eccezione dichiarata dei prospetti di costo, dove le voci di costo sono positive perché si sommano; e l'indice temporale $t=0$ indica il momento del rogito, non il primo anno di gestione.

\section{Capitalizzazione, sconto, e la convenzione che cambia il risultato}

\subsection{Le due operazioni elementari}

\begin{definizione}[Montante]
Un capitale $X$ investito per $m$ periodi a tasso di periodo $j$ costante, con reinvestimento degli interessi, produce il montante
\begin{equation}
X_m = X\,(1+j)^m .
\end{equation}
\end{definizione}

La formula si ottiene per induzione dal solo fatto che gli interessi di un periodo entrano nella base del successivo: $X_1 = X(1+j)$, e se $X_m = X(1+j)^m$ allora $X_{m+1} = X_m(1+j) = X(1+j)^{m+1}$. Il regime alternativo, quello semplice, in cui gli interessi non entrano nella base e $X_m = X(1+jm)$, non compare in questo modello: nessuno dei contratti trattati lo usa.

\begin{definizione}[Valore attuale]
Il valore oggi di un importo $Y$ disponibile fra $m$ periodi, al tasso di sconto di periodo $j$, è
\begin{equation}
\VAN(Y) = \frac{Y}{(1+j)^m} .
\end{equation}
\end{definizione}

Le due operazioni sono l'una l'inversa dell'altra, e questo è tutto ciò che serve: ogni formula che segue è una combinazione lineare di valori attuali.

\subsection{Tassi equivalenti, e perché la convenzione italiana non è un dettaglio}

Un tasso annuo $i$ si converte in un tasso mensile in due modi diversi, che danno numeri diversi.

\begin{equation}
j_{\text{lineare}} = \frac{i}{12}, \qquad
j_{\text{composto}} = (1+i)^{1/12} - 1 .
\label{eq:conversione}
\end{equation}

\paragraph{In parole.} La prima riga dice: per ottenere il tasso di un mese, dividi il tasso dell'anno per dodici. Su un tasso annuo del $3{,}2$ per cento fa $0{,}032/12 = 0{,}0026667$, cioe' circa un quarto di punto al mese. La seconda dice: prendi il numero $1{,}032$, estraine la radice dodicesima, e togli uno. Fa $0{,}0026281$. Le due rispondono alla stessa domanda e danno numeri diversi perche' la prima ignora che gli interessi del primo mese producono interessi negli undici successivi. La banca italiana usa la prima, quindi il modello usa la prima, ma la mostra come cella modificabile perche' chi confronta un mutuo italiano con un prodotto estero deve sapere che le due convenzioni non sono la stessa cosa.

Il secondo è il tasso \emph{finanziariamente equivalente}, cioè quello che capitalizzato dodici volte restituisce esattamente $i$: $(1+j_{\text{composto}})^{12} = 1+i$. Il primo non ha questa proprietà, perché $(1+i/12)^{12} > 1+i$ per $i>0$: dividere per dodici e poi capitalizzare dodici volte produce più di quanto si è dichiarato.

\begin{osservazione}
Su $i = 0{,}032$ si ha $j_{\text{lineare}} = 0{,}0026667$ e $j_{\text{composto}} = 0{,}0026281$, cioè una differenza di circa il quattro per cento sul tasso di periodo. Capitalizzata su venticinque anni, la differenza fra i due regimi sposta gli interessi totali di alcune centinaia di euro. La convenzione dei contratti di mutuo italiani è la prima, cioè la divisione lineare, e il modello la adotta come predefinita perché deve riprodurre il piano che la banca calcola, non quello che sarebbe finanziariamente coerente. Il foglio \texttt{Simulatore mutuo} espone la scelta come cella, così che la differenza si veda invece di essere subita.
\end{osservazione}

Questo è il primo esempio di un criterio che ritorna in tutto il documento: fra due formule algebricamente vicine si sceglie quella che riproduce il contratto, non quella più elegante, e la scelta si dichiara.

\section{Rendite: il valore di una successione di pagamenti}

Il modello valuta successioni di pagamenti in tre forme: costanti, crescenti a tasso costante, e arbitrarie. Le prime due hanno una forma chiusa che vale derivare, perché il workbook la usa dentro celle singole.

\subsection{Rendita costante posticipata}

\begin{proposizione}[Fattore di annualità]
Il valore attuale di $n$ pagamenti unitari, il primo alla fine del primo periodo, scontati al tasso di periodo $s>0$, vale
\begin{equation}
a_{n,s} \;=\; \sum_{k=1}^{n} \frac{1}{(1+s)^k} \;=\; \frac{1-(1+s)^{-n}}{s} .
\label{eq:annualita}
\end{equation}

\paragraph{In parole.} A sinistra c'e' la domanda: quanto vale oggi ricevere un euro alla fine di ognuno dei prossimi $n$ anni. A destra c'e' la risposta in forma compatta. La somma dice di attualizzare ogni euro per il numero di anni che lo separa da oggi e poi sommare tutto; la forma chiusa dice che quella somma si calcola in un colpo, senza scrivere gli $n$ addendi.

Con $n = 25$ anni e $s = 3$ per cento il fattore vale $17{,}41$: venticinque euro incassati uno all'anno per venticinque anni valgono oggi diciassette euro e quaranta, non venticinque. La differenza di sette euro e sessanta e' il prezzo dell'attesa, e cresce con il tasso di sconto: al sei per cento lo stesso flusso vale $12{,}78$.
\end{proposizione}

\begin{proof}
Si ponga $v = (1+s)^{-1}$, quindi $0<v<1$. La somma è geometrica di primo termine $v$ e ragione $v$:
\[
\sum_{k=1}^{n} v^k = v\,\frac{1-v^n}{1-v} .
\]
Sostituendo $1-v = 1-\frac{1}{1+s} = \frac{s}{1+s}$ si ottiene
\[
v\,\frac{1-v^n}{1-v} = \frac{1}{1+s}\cdot\frac{1-v^n}{s/(1+s)} = \frac{1-v^n}{s} = \frac{1-(1+s)^{-n}}{s},
\]
che è la tesi. Il caso $s=0$ va trattato a parte e dà $a_{n,0}=n$, come si vede direttamente dalla somma.
\end{proof}

\subsection{Rendita crescente a tasso costante}

È la forma che serve per un canone indicizzato, e la sua derivazione è quella che il modello implementa in una cella.

\begin{proposizione}[Fattore di una rendita crescente]
Il valore attuale di $n$ pagamenti, il primo unitario e ciascuno successivo maggiore del precedente del fattore $(1+g)$, scontati al tasso $s$, vale
\begin{equation}
F(g,s,n) \;=\; \sum_{k=1}^{n} \frac{(1+g)^{k-1}}{(1+s)^k}
\;=\;
\begin{cases}
\dfrac{q\,(1-q^{\,n})}{(1+g)\,(1-q)}, & q \neq 1,\\[1.2em]
\dfrac{n}{1+g}, & q = 1,
\end{cases}
\qquad q := \frac{1+g}{1+s}.
\label{eq:crescente}
\end{equation}

\paragraph{In parole.} E' la formula precedente con una complicazione: i pagamenti non sono tutti uguali, ognuno e' piu' grande del precedente di una percentuale fissa $g$. E' il caso di un canone di affitto indicizzato: il primo anno incassi cento, il secondo cento piu' l'indicizzazione, e cosi' via.

La lettera $q$, definita a destra della formula, e' il rapporto fra quanto il pagamento cresce e quanto lo si sconta. Se $q$ e' minore di uno lo sconto vince sulla crescita e la somma converge a un numero finito anche su tanti anni; se $q$ e' maggiore di uno la crescita vince e la somma diventa grande. Il caso $q = 1$, cioe' crescita esattamente pari al tasso di sconto, e' scritto a parte perche' li' la formula generale dividerebbe per zero: quando le due cose si annullano, ogni pagamento vale oggi la stessa cifra e la somma e' semplicemente quella cifra moltiplicata per il numero di anni.

Sul caso del progetto, con inflazione al due per cento e sconto al tre, il fattore vale $21{,}64$: un canone indicizzato valeva $21{,}64$ euro per ogni euro di canone iniziale, contro i $17{,}41$ di un canone fermo. La differenza fra i due numeri e' il valore dell'indicizzazione, ed e' il calcolo del capitolo sull'inflazione.
\end{proposizione}

\begin{proof}
Si raccoglie $(1+g)^{-1}$ per allineare gli esponenti:
\[
\sum_{k=1}^{n} \frac{(1+g)^{k-1}}{(1+s)^k}
= \frac{1}{1+g}\sum_{k=1}^{n} \frac{(1+g)^{k}}{(1+s)^k}
= \frac{1}{1+g}\sum_{k=1}^{n} q^{k}.
\]
La somma geometrica di ragione $q$ vale $q(1-q^n)/(1-q)$ per $q\neq 1$, da cui la prima riga. Se $q=1$, cioè $g=s$, ogni addendo $q^k$ vale $1$ e la somma è $n$, da cui la seconda.
\end{proof}

\begin{osservazione}[Perché la forma chiusa, e perché il caso singolare]
La somma esplicita è più leggibile e il motore Python la calcola così. Il workbook non può: una somma su $n$ termini con $n$ variabile richiede una formula matriciale oppure una colonna di appoggio, e nessuna delle due è ispezionabile come una cella singola. Il foglio usa quindi la \eqref{eq:crescente}, con il ramo $q=1$ selezionato da un confronto a tolleranza. Il caso singolare non è di scuola: assumere una crescita del canone pari al tasso di sconto è un'ipotesi che si incontra, e senza il ramo dedicato la cella darebbe una divisione per zero. Un test del progetto verifica che forma chiusa e somma esplicita coincidano su sei configurazioni, incluso $g=s$: è la doppia implementazione applicata a una formula invece che a un modello.
\end{osservazione}

\subsection{Serie di flussi arbitrari}

Quando i pagamenti non hanno una legge, il valore attuale è la somma diretta
\begin{equation}
\VAN(s) \;=\; \sum_{t=0}^{N} \frac{F_t}{(1+s)^t},
\label{eq:van}
\end{equation}

\paragraph{In parole.} Prendi ogni flusso di cassa, dividilo per uno piu' il tasso di sconto elevato al numero di anni che lo separa da oggi, e somma tutto. Il flusso del tempo zero non si divide per niente, perche' e' gia' oggi: nella formula questo si vede dal fatto che con $t=0$ il denominatore vale $(1+s)^0 = 1$.

Il risultato risponde a una domanda sola: mettendo tutto in euro di oggi, questa operazione mi lascia piu' o meno di quanto ci ho messo. Se il numero e' positivo l'operazione batte il tasso di sconto scelto, se e' negativo non lo batte.
con $F_0$ tipicamente negativo, perché è l'esborso iniziale. Si noti che l'indice parte da zero e che $F_0$ non viene scontato: è già in euro di oggi. È una banalità che produce un errore ricorrente, cioè scontare anche il primo termine, con l'effetto di sottostimare l'esborso e sovrastimare il valore attuale netto.

\section{Inflazione: dal nominale al reale}

\subsection{L'equazione di Fisher, e l'errore della sottrazione}

\begin{definizione}[Rendimento reale]
Se un investimento rende $r$ in termini nominali su un periodo in cui il livello dei prezzi cresce di $\pi$, il rendimento reale $r^*$ è definito dal rapporto fra il potere d'acquisto finale e quello iniziale:
\begin{equation}
1 + r^* \;=\; \frac{1+r}{1+\pi}
\qquad\Longleftrightarrow\qquad
r^* \;=\; \frac{1+r}{1+\pi} - 1 \;=\; \frac{r-\pi}{1+\pi}.
\label{eq:fisher}
\end{equation}

\paragraph{In parole.} La prima uguaglianza e' la definizione e va letta cosi': un euro investito diventa $(1+r)$ euro nominali, ma nel frattempo i prezzi sono cresciuti di $(1+\pi)$, quindi quei $(1+r)$ euro comprano $(1+r)/(1+\pi)$ volte quello che comprava un euro all'inizio. Quel rapporto, meno uno, e' il rendimento in potere d'acquisto.

Con un rendimento nominale del cinque per cento e un'inflazione del due: cento euro diventano centocinque, ma per comprare quello che prima costava cento adesso servono centodue, quindi in termini di cose comprabili si e' guadagnato $105/102 = 1{,}0294$, cioe' il $2{,}94$ per cento e non il tre.
\end{definizione}

La forma che si incontra quasi sempre, $r^* \approx r - \pi$, è l'approssimazione al primo ordine della \eqref{eq:fisher}. L'errore che si commette usandola è esattamente
\begin{equation}
(r-\pi) - r^* \;=\; (r-\pi) - \frac{r-\pi}{1+\pi} \;=\; (r-\pi)\,\frac{\pi}{1+\pi},
\label{eq:errore-fisher}
\end{equation}

\paragraph{In parole.} La formula misura di quanto sbaglia chi usa la sottrazione. Si legge: l'errore e' lo scarto approssimato, moltiplicato per l'inflazione, diviso uno piu' l'inflazione. Il fatto che l'inflazione compaia come fattore dice la cosa utile: con inflazione nulla l'errore e' nullo e la sottrazione e' esatta, e piu' l'inflazione sale piu' la sottrazione inganna. Con inflazione al due per cento l'errore e' circa il due per cento dello scarto, con inflazione al dieci diventa circa il nove.
quindi è proporzionale al prodotto fra il rendimento reale approssimato e l'inflazione. È piccolo, non è nullo, ed è sempre dello stesso segno quando $r>\pi>0$: la sottrazione sovrastima il rendimento reale.

\begin{osservazione}
Con $r=0{,}05$ e $\pi=0{,}02$ si ha $r^*=0{,}029412$ contro $0{,}03$ della sottrazione: sei centesimi di punto, cioè il due per cento del rendimento reale stesso. Su un orizzonte lungo l'errore si compone: $1{,}03^{25}/1{,}029412^{25} = 1{,}0148$, cioè un uno e mezzo per cento di montante finale attribuito a una scelta di formula. Il modello usa sempre la \eqref{eq:fisher} e mostra accanto l'errore \eqref{eq:errore-fisher}, perché la sottrazione resta utile per il conto a mente e va saputa per quello che è.
\end{osservazione}

\subsection{Deflazione di un importo futuro}

Un importo nominale $Y$ disponibile fra $N$ anni vale, in potere d'acquisto di oggi,
\begin{equation}
Y^* = \frac{Y}{(1+\pi)^N}.
\label{eq:deflaziona}
\end{equation}
È la stessa operazione dello sconto, con l'inflazione al posto del tasso di sconto, e risponde a una domanda diversa: non quanto vale oggi quel denaro per me che potrei investirlo, ma quante cose comprerà quando lo avrò.

\section{Il costo dell'operazione}

\subsection{La base imponibile: una funzione a tratti}

L'imposta di registro sul trasferimento non si applica al prezzo ma a una base imponibile che dipende da un'opzione dell'acquirente. Esercitando l'opzione prezzo-valore, ammessa per le persone fisiche che acquistano immobili a uso abitativo, la base diventa il valore catastale
\begin{equation}
V_c \;=\; R \cdot \kappa \cdot \mu(\alpha),
\qquad
\mu(\alpha) = \begin{cases}\mu_1 & \text{con agevolazione},\\ \mu_0 & \text{senza},\end{cases}
\label{eq:catastale}
\end{equation}
dove $\kappa$ è il coefficiente di rivalutazione della rendita e $\mu$ il moltiplicatore di categoria. La base imponibile è quindi
\begin{equation}
B(P) \;=\;
\begin{cases}
V_c, & \text{se si opta per il prezzo-valore e } R>0,\\
P, & \text{altrimenti.}
\end{cases}
\label{eq:base}
\end{equation}

\paragraph{In parole.} La base su cui si calcola l'imposta di registro non e' sempre il prezzo. Se l'acquirente e' una persona fisica, l'immobile e' abitativo e si chiede espressamente in atto l'opzione prezzo-valore, allora la base diventa il valore catastale, che si ricava dalla rendita e che di norma e' molto piu' basso del prezzo. Altrimenti la base e' il prezzo.

Sul caso del progetto: prezzo centoventimila, valore catastale $51\,975$. L'imposta al due per cento sul secondo fa $1\,039$ euro, sul primo ne farebbe $2\,400$. La differenza e' il valore dell'opzione, e si esercita chiedendola al notaio: non e' automatica.

Il punto cruciale, che ritorna nel capitolo sull'inversione, è che nel primo ramo $B$ \emph{non dipende da $P$}: l'imposta di registro diventa una costante rispetto al prezzo. È la leva fiscale più grossa dell'operazione, e la sua natura matematica è quella di rendere costante un termine che altrimenti sarebbe proporzionale.

\subsection{Le imposte di trasferimento}

Nei due regimi si ha
\begin{equation}
T_a(P) \;=\;
\begin{cases}
\max\big(B(P)\,\tau(\alpha),\; T_{\min}\big) + \iota + \gamma, & \text{venditore privato o esente,}\\[0.4em]
P\,\upsilon(\alpha,\lambda) + 3\,\phi, & \text{venditore impresa con IVA,}
\end{cases}
\label{eq:imposte}
\end{equation}

\paragraph{In parole.} Le due righe sono i due mondi in cui puo' cadere un acquisto, e non sono varianti l'una dell'altra.

Prima riga, venditore privato. Si applica l'aliquota di registro alla base, ma se il risultato e' inferiore al minimo di legge si paga il minimo: e' cio' che dice il $\max$. A questo si sommano due imposte in misura fissa, ipotecaria e catastale, che non dipendono dal prezzo.

Seconda riga, venditore impresa con IVA. Si applica l'aliquota IVA al prezzo intero, e si sommano tre imposte fisse invece di due. Nessun $\max$, perche' non c'e' un minimo di legge, e nessun ancoraggio al catastale, perche' il prezzo-valore non si applica al regime IVA.

Il salto fra le due righe, su ottantamila euro con agevolazione, e' da circa millecinquecento a circa tremilacinquecento euro.
dove $\tau$ è l'aliquota di registro, $T_{\min}$ il minimo di legge, $\iota$ e $\gamma$ le imposte ipotecaria e catastale in misura fissa, $\upsilon$ l'aliquota IVA che dipende dall'agevolazione e dalla qualifica di lusso $\lambda$, e $\phi$ l'imposta fissa che nel regime IVA si applica tre volte.

Due non linearità vanno segnalate perché rompono la linearità che il capitolo sull'inversione sfrutta. La prima è il $\max$ con il minimo di legge, che rende $T_a$ costante per prezzi sufficientemente bassi. La seconda è la dipendenza di $\upsilon$ da $\lambda$, che è una variabile discreta e non un parametro continuo.

\begin{osservazione}[Il salto fra i due regimi]
Le due righe della \eqref{eq:imposte} non sono perturbazioni l'una dell'altra: su un prezzo di ottantamila euro con agevolazione, il regime privato con prezzo-valore produce circa millecinquecento euro di imposta, quello IVA circa tremilacinquecento. Il rapporto è maggiore di due e cresce col prezzo, perché l'IVA colpisce il prezzo intero mentre il registro con prezzo-valore resta ancorato alla rendita. È la ragione per cui una graduatoria che applichi lo stesso regime a un usato da privato e a un nuovo da costruttore non è imprecisa: è ordinata al contrario nella parte alta.
\end{osservazione}

\subsection{Il costo totale come funzione affine del prezzo}

I costi accessori si compongono di imposte, provvigione, notaio, altri costi e oneri del mutuo:
\begin{equation}
A(P) \;=\; T_a(P) \;+\; P\,p\,(1+\upsilon_p) \;+\; \nu \;+\; \omega \;+\; \Omega_M,
\label{eq:accessori}
\end{equation}
con $p$ aliquota di provvigione, $\upsilon_p$ IVA sulla provvigione, $\nu$ onorario notarile, $\omega$ altri costi, $\Omega_M$ oneri del mutuo. Il costo totale è $C(P) = P + A(P)$.

\begin{proposizione}[Forma affine]
Fuori dal tratto in cui il minimo di legge è vincolante, il costo totale è affine nel prezzo:
\begin{equation}
C(P) \;=\; P\,(1+k) + c,
\label{eq:costo-affine}
\end{equation}

\paragraph{In parole.} Il costo totale dell'operazione si scrive come il prezzo moltiplicato per un fattore, piu' una cifra fissa. Il fattore $(1+k)$ dice quanti euro escono dal conto per ogni euro di prezzo: uno e' il prezzo stesso, $k$ sono i costi che crescono col prezzo. La cifra fissa $c$ e' tutto quello che si paga comunque, che il prezzo sia alto o basso.

Sul caso del progetto $k$ vale il $3{,}7$ per cento e $c$ vale $7\,165$ euro. Quindi comprando a centoventimila si spende $120\,000 \times 1{,}037 + 7\,165 = 131\,600$ circa, e comprando a sessantamila si spende $62\,200 + 7\,165 = 69\,400$ circa. Si noti la conseguenza: dimezzando il prezzo il costo totale non si dimezza, perche' i $7\,165$ euro restano. E' l'origine dell'errore che il capitolo sull'inversione corregge.
con i coefficienti
\begin{align}
k &= \underbrace{\begin{cases}
\upsilon(\alpha,\lambda) & \text{regime IVA},\\
0 & \text{registro con prezzo-valore},\\
\tau(\alpha) & \text{registro senza prezzo-valore},
\end{cases}}_{\text{quota fiscale marginale}}
\;+\; p\,(1+\upsilon_p),
\label{eq:k}\\[0.6em]
c &= \begin{cases}
3\phi & \text{regime IVA},\\
\max(V_c\,\tau(\alpha), T_{\min}) + \iota + \gamma & \text{registro con prezzo-valore},\\
\iota + \gamma & \text{registro senza prezzo-valore},
\end{cases}
\;+\; \nu + \omega + \Omega_M .
\label{eq:c}
\end{align}

\paragraph{In parole.} Le due formule dicono da dove vengono i due numeri della formula precedente, caso per caso.

Il primo, $k$, e' quanto di ogni euro di prezzo si trasforma in costo aggiuntivo: nel regime IVA e' l'aliquota IVA, con il prezzo-valore attivo e' zero perche' il registro non segue il prezzo, senza prezzo-valore e' l'aliquota di registro. A questo si somma sempre la provvigione con la sua IVA, che e' una percentuale del prezzo in ogni regime.

Il secondo, $c$, e' tutto quello che non dipende dal prezzo: le imposte fisse, l'onorario del notaio, gli altri costi, gli oneri del mutuo, e con il prezzo-valore anche l'intera imposta di registro, che essendo ancorata alla rendita catastale e' una costante rispetto al prezzo.

La riga da guardare e' la seconda del caso $c$: e' controintuitivo che un'imposta proporzionale finisca fra i costi fissi, e succede solo perche' la proporzionalita' e' rispetto al catastale e non al prezzo.
\end{proposizione}

La proposizione è la sostanza del capitolo, e la sua conseguenza è controintuitiva: con il prezzo-valore attivo l'intera imposta di registro finisce in $c$, cioè nella parte \emph{fissa}, e non in $k$. L'incidenza percentuale dei costi accessori, $A(P)/P$, non è quindi una costante del modello ma una funzione decrescente del prezzo. Trattarla come costante è l'errore che il capitolo~\ref{sec:inversione} corregge.

\subsection{Il denominatore dei rendimenti}

\begin{osservazione}
Ogni rendimento è un rapporto, e la scelta del denominatore è una scelta di modello. Usare $P$ invece di $C$ ignora imposte, notaio, provvigione e oneri del mutuo, che sono capitale uscito dal conto e che alla rivendita non torna. La distorsione è sistematica, sempre nello stesso verso, e vale $C/P - 1$, cioè l'incidenza dei costi accessori: sul caso di riferimento del progetto il $9{,}63$ per cento. Un modello che sbaglia sempre nella stessa direzione è peggio di uno rumoroso, perché sembra affidabile. Il modello usa quindi $C$ per i rendimenti dell'attivo e $E$, il capitale proprio, per il solo \emph{cash on cash}, che misura l'operazione finanziata e non l'immobile.
\end{osservazione}

\section{Il mutuo: ammortamento alla francese}

\subsection{Derivazione della rata}

\begin{proposizione}[Rata costante]
Un capitale $M$ da restituire in $n$ rate posticipate costanti, al tasso di periodo $j$, richiede la rata
\begin{equation}
\Gamma \;=\; M\,\frac{j}{1-(1+j)^{-n}} \;=\; \frac{M}{a_{n,j}} .
\label{eq:rata}
\end{equation}

\paragraph{In parole.} La rata si ottiene dividendo il capitale per il fattore di annualita' della formula \eqref{eq:annualita}. Il significato e' questo: il fattore dice quanto vale oggi un euro di rata ripetuto per tutta la durata, quindi per sapere quanta rata serve a coprire il capitale basta dividere il capitale per quel valore.

Su novantamila euro al $3{,}2$ per cento per venticinque anni, cioe' trecento rate: il tasso mensile e' $0{,}032/12 = 0{,}00266\overline{6}$, il fattore di annualita' su trecento mesi vale $206{,}3$, e la rata e' $90\,000/206{,}3 = 436{,}21$ euro. Il totale delle rate e' $436{,}21 \times 300 = 130\,863$ euro, quindi gli interessi complessivi sono $40\,863$ euro: quasi la metà del capitale prestato.
\end{proposizione}

\begin{proof}
La condizione che definisce il piano è l'equivalenza finanziaria al tempo zero fra l'erogato e la successione delle rate: il valore attuale delle rate, scontate al tasso del contratto, deve essere pari al capitale ricevuto. Quindi
\[
M \;=\; \sum_{k=1}^{n} \frac{\Gamma}{(1+j)^k} \;=\; \Gamma\,a_{n,j},
\]
dove l'ultima uguaglianza è la \eqref{eq:annualita}. Risolvendo per $\Gamma$ e sostituendo la forma chiusa di $a_{n,j}$ si ottiene la tesi.
\end{proof}

È la funzione che i fogli di calcolo chiamano \texttt{PMT} o \texttt{RATA}. Vale notare che la \eqref{eq:rata} non contiene alcuna ipotesi sulla composizione fra quota capitale e quota interessi: quella composizione è una \emph{conseguenza} della rata costante, non una scelta.

\subsection{Debito residuo, quota interessi, quota capitale}

\begin{proposizione}[Debito residuo in forma chiusa]
Dopo la rata $k$ il debito residuo vale
\begin{equation}
D_k \;=\; M\,\frac{(1+j)^n - (1+j)^k}{(1+j)^n - 1} .
\label{eq:residuo}
\end{equation}

\paragraph{In parole.} La formula dice quanto si deve ancora alla banca dopo aver pagato $k$ rate, senza dover percorrere tutto il piano una rata per volta.

Il modo di leggerla e' guardare i due estremi. Con $k=0$, cioe' prima di pagare qualsiasi cosa, il numeratore diventa $(1+j)^n - 1$, uguale al denominatore, quindi il debito residuo e' esattamente $M$, il capitale erogato. Con $k=n$, cioe' dopo l'ultima rata, il numeratore diventa zero e il debito residuo e' zero. In mezzo la formula interpola, e non lo fa in modo lineare: dopo aver pagato metà delle rate il debito residuo non e' metà del capitale ma di piu', perche' le prime rate sono composte in prevalenza di interessi.

Sul caso del progetto, a metà del piano, cioe' dopo centocinquanta rate su trecento, restano circa $52\,000$ euro su novantamila: si e' pagato oltre la metà delle rate e restituito meno del quaranta per cento del capitale.
\end{proposizione}

\begin{proof}
La ricorsione del piano è $D_k = D_{k-1}(1+j) - \Gamma$ con $D_0 = M$. Si tratta di una ricorsione affine, la cui soluzione si ottiene cercando il punto fisso $D^\star$ tale che $D^\star = D^\star(1+j)-\Gamma$, cioè $D^\star = \Gamma/j$, e osservando che lo scarto dal punto fisso si moltiplica per $(1+j)$ a ogni passo:
\[
D_k - \frac{\Gamma}{j} \;=\; (1+j)^k\Big(M - \frac{\Gamma}{j}\Big)
\quad\Longrightarrow\quad
D_k \;=\; \frac{\Gamma}{j} + (1+j)^k\Big(M-\frac{\Gamma}{j}\Big).
\]
Sostituendo $\Gamma = Mj/(1-(1+j)^{-n})$ si ha $\Gamma/j = M/(1-(1+j)^{-n})$, e con qualche passaggio algebrico
\[
D_k = \frac{M}{1-(1+j)^{-n}}\Big(1-(1+j)^{k-n}\Big)
= M\,\frac{(1+j)^n-(1+j)^k}{(1+j)^n-1},
\]
che è la tesi. Si verifica immediatamente che $D_0 = M$ e $D_n = 0$.
\end{proof}

Da qui seguono le due scomposizioni della rata:
\begin{equation}
I_k \;=\; j\,D_{k-1}, \qquad K_k \;=\; \Gamma - I_k,
\label{eq:scomposizione}
\end{equation}

\paragraph{In parole.} Ogni rata si divide in due pezzi. Il primo, gli interessi, si calcola applicando il tasso del mese al debito che c'era all'inizio del mese: e' l'affitto del denaro per quel mese. Il secondo, la quota capitale, e' quello che resta della rata dopo aver pagato gli interessi, ed e' l'unica parte che riduce il debito.

Sulla prima rata del caso del progetto: il debito e' novantamila, il tasso mensile e' $0{,}0026\overline{6}$, quindi gli interessi sono $240$ euro e la quota capitale e' $436{,}21 - 240 = 196{,}21$. Sull'ultima rata gli interessi sono circa un euro e la quota capitale e' quasi tutta la rata. Poiche' il debito scende, gli interessi scendono, e poiche' la rata e' fissa la quota capitale sale: e' un'aritmetica automatica e non una scelta della banca.
cioè la quota interessi è il tasso applicato al debito \emph{di inizio periodo} e la quota capitale è il residuo. La forma chiusa \eqref{eq:residuo} è ciò che permette al modello di calcolare il debito residuo a un'epoca qualunque senza percorrere il piano, e il foglio degli scenari la usa per restare esatto quando si cambia il tasso di scenario.

\begin{osservazione}[Perché all'inizio si pagano soprattutto interessi]
Dalla \eqref{eq:scomposizione} si vede che $I_k$ è decrescente perché $D_{k-1}$ lo è, e quindi $K_k$ è crescente. Non è una scelta della banca né una clausola: è una conseguenza aritmetica della rata costante. La conseguenza pratica, che il modello dichiara nelle note, è però l'opposto di quella che si sente ripetere: il fatto che all'inizio si paghino soprattutto interessi \emph{non} è una ragione per rimborsare presto, perché ogni mese la scelta è la stessa, cioè estinguere ora oppure pagare $j$ sul residuo per rimandare la decisione. La condizione di convenienza del rimborso è che nessun impiego alternativo renda, al netto dell'imposta, almeno $i$.
\end{osservazione}

\subsection{Rimborsi volontari e le due modalità di imputazione}

Con un versamento aggiuntivo $x_k$ al periodo $k$, la ricorsione diventa
\begin{equation}
D_k \;=\; \max\big(0,\; D_{k-1}(1+j_k) - \Gamma_k - x_k\big),
\label{eq:ricorsione}
\end{equation}
e la definizione di $\Gamma_k$ dipende dalla modalità dichiarata alla banca:
\begin{equation}
\Gamma_k \;=\;
\begin{cases}
\Gamma, & \text{riduzione della durata,}\\[0.4em]
D_{k-1}\,\dfrac{j_k}{1-(1+j_k)^{-(n-k+1)}}, & \text{riduzione della rata.}
\end{cases}
\label{eq:modalita}
\end{equation}

Nel primo caso la rata resta quella originaria e il piano si accorcia, perché $D_k$ raggiunge zero prima di $n$; nel secondo la scadenza resta fissa e la rata si ricalcola sul residuo e sui periodi rimanenti. Le due producono risultati molto diversi: sul caso di riferimento, con cento euro al mese di versamento aggiuntivo, la riduzione della durata risparmia $11\,373$ euro di interessi contro $6\,543$ della riduzione della rata. La differenza non è un dettaglio contrattuale: è la scelta fra ridurre l'esposizione totale al tasso e ridurre l'impegno mensile, e va dichiarata perché non la sceglie il foglio.

\subsection{Percorso del tasso a gradini}

Su un mutuo indicizzato il tasso non è costante. Il modello descrive il percorso come una funzione a gradini: dati $m$ punti $(\hat{k}_1,\delta_1),\dots,(\hat{k}_m,\delta_m)$ con $\hat{k}_1<\dots<\hat{k}_m$, dove $\delta_h$ è la variazione cumulata rispetto al tasso iniziale, il tasso del periodo $k$ è
\begin{equation}
i_k \;=\; i \;+\; \delta_{h(k)},
\qquad
h(k) = \max\{\,h : \hat{k}_h \le k\,\},
\label{eq:gradini}
\end{equation}

\paragraph{In parole.} Il tasso del mese $k$ e' il tasso iniziale piu' la variazione del gradino piu' avanzato fra quelli il cui mese e' gia' arrivato. La scrittura con il $\max$ dell'insieme dice esattamente questo: fra tutti i gradini il cui mese di partenza e' minore o uguale al mese corrente, prendi quello con il numero piu' alto.

Con tre gradini a piu' $1{,}26$ dal mese tredici, piu' $2{,}52$ dal diciassette e piu' $3{,}78$ dal ventuno, e un tasso iniziale del $3{,}2$ per cento: nei primi dodici mesi si paga il $3{,}2$, dal tredicesimo il $4{,}46$, dal diciassettesimo il $5{,}72$, dal ventunesimo il $6{,}98$ e cosi' resta. Le variazioni sono cumulate e non si sommano fra loro: il terzo gradino dice piu' $3{,}78$ rispetto all'inizio, non piu' $3{,}78$ rispetto al secondo.
con la convenzione $\delta_{h}=0$ se nessun gradino è stato raggiunto. La \eqref{eq:gradini} entra nella \eqref{eq:ricorsione} attraverso $j_k = i_k/12$, e la ricorsione va percorsa numericamente perché non esiste forma chiusa con tasso variabile.

\begin{osservazione}[Il piano che non si chiude]
Con la modalità a durata variabile e un rialzo forte, la \eqref{eq:ricorsione} può non raggiungere lo zero entro l'orizzonte tabulato. Il modello tabula $n_{\max}=480$ periodi, cioè quarant'anni: se $D_{480}>0$ il piano non è risolto ma troncato, e le grandezze di sintesi calcolate sulla tabella, cioè durata effettiva e interessi totali, descrivono soltanto la porzione tabulata. Il foglio lo dichiara con una riga di esito dedicata, perché due numeri veri che rispondono a una domanda diversa da quella posta sono peggio di un errore.
\end{osservazione}

\subsection{Costo effettivo del finanziamento}

Il costo effettivo non è il tasso nominale, perché il nominale ignora gli oneri accessori e le eventuali variazioni. Si definisce come il tasso interno della successione dei flussi del solo mutuo:
\begin{equation}
\text{trovare } y \text{ tale che} \quad
M - \Omega_M \;-\; \sum_{k=1}^{n} \frac{\Gamma_k + x_k}{(1+y)^k} \;=\; 0,
\label{eq:costo-effettivo}
\end{equation}
dove $M-\Omega_M$ è l'incasso netto al tempo zero, perché imposta sostitutiva, istruttoria, perizia e notaio del mutuo si pagano subito. Il tasso annuo effettivo è $12y$ nella convenzione lineare del modello.

\begin{osservazione}
La \eqref{eq:costo-effettivo} non è il TAEG di legge, che segue la convenzione della Banca d'Italia sull'insieme delle voci da includere. Il modello lo dichiara: è una stima interna, utile per confrontare due preventivi sulla stessa base, non un valore da citare in una contestazione.
\end{osservazione}

\section{Il reddito da locazione}

\subsection{Dalla potenzialità al ricavo effettivo}

Il canone contrattuale è una potenzialità, non un incasso. Con $L$ canone annuo a pieno regime, $\sigma$ mesi di sfitto attesi all'anno e $\beta$ accantonamento per morosità come quota del canone,
\begin{equation}
L_e \;=\; \Big(L - \frac{L}{12}\,\sigma\Big)\,(1-\beta)
\;=\; L\,\Big(1-\frac{\sigma}{12}\Big)(1-\beta).
\label{eq:ricavo-effettivo}
\end{equation}

\paragraph{In parole.} Il canone che entra davvero e' meno di quello scritto sul contratto, per due ragioni che si applicano una dopo l'altra. Prima si toglie la parte dell'anno in cui l'immobile e' vuoto: se si attende un mese di sfitto, si incassano undici dodicesimi. Poi, su quello che resta, si mette da parte una quota per il rischio che l'inquilino non paghi.

Su un canone di seicento euro al mese, cioe' $7\,200$ l'anno, con un mese di sfitto e il tre per cento di morosita': $7\,200 \times (11/12) = 6\,600$, e $6\,600 \times 0{,}97 = 6\,402$. Il canone dichiarato era $7\,200$, quello su cui contare e' $6\,402$, cioe' l'undici per cento in meno.

L'ordine conta: applicare la morosita' dopo lo sfitto e' prudenziale, perche' si accantona sul minore dei due importi. Fare il contrario darebbe un numero leggermente piu' alto.
Le due correzioni sono moltiplicative e non additive, cioè la morosità si applica a ciò che resta dopo lo sfitto: è la scelta prudenziale, e la differenza rispetto alla forma additiva è di secondo ordine ma sempre dello stesso segno.

\begin{osservazione}
$\sigma$ e $\beta$ sono i due parametri che chi valuta tende a mettere a zero, ed è l'errore più comune nelle valutazioni ottimistiche. Un mese l'anno di sfitto è l'assunzione prudenziale standard, non pessimistica, perché fra un contratto e l'altro il vuoto reale è di norma più lungo; e uno sfratto per morosità richiede circa un anno e mezzo, durante il quale il canone non entra e le imposte si pagano comunque.
\end{osservazione}

\subsection{Reddito operativo netto e utile netto}

I costi del proprietario, tutti annui, sono le spese condominiali a suo carico $\chi$, la manutenzione ordinaria come quota $\mathrm{m}$ del valore, l'assicurazione $\zeta$, il costo figurativo del tempo $\eta$, l'accantonamento per la ristrutturazione di fine ciclo, e l'IMU. Il reddito operativo netto è
\begin{equation}
\mathrm{NOI} \;=\; L_e \;-\; \Big(\chi + P\,\mathrm{m} + \zeta + \eta + \frac{P\,\mathrm{r}}{N_r} + \mathrm{IMU}\Big),
\label{eq:noi}
\end{equation}

\paragraph{In parole.} Dal canone effettivo si toglie tutto quello che l'immobile costa al proprietario in un anno, e resta il reddito operativo netto. Le voci in parentesi, nell'ordine: le spese condominiali che restano a carico suo, la manutenzione ordinaria calcolata come quota del valore, l'assicurazione del fabbricato, il valore del proprio tempo se lo si vuole contare, l'accantonamento annuo per il rifacimento di fine ciclo, e l'IMU.

Le due voci che si dimenticano piu' spesso sono la quinta e la prima. La quinta, l'accantonamento, e' l'unica riga che rappresenta un costo che non si paga quest'anno: un immobile che si tiene quarant'anni va rifatto almeno una volta, e se il rifacimento costa un terzo del valore, allora ogni anno se ne consuma un quarantesimo. Ignorarla fa apparire un rendimento che non esiste. La prima, il condominio, si dimentica perche' non e' una voce dell'annuncio.

Sul caso del progetto: canone effettivo $5\,335$, costi $3\,530$, reddito operativo netto $1\,805$ euro.
dove $\mathrm{r}$ è il costo di un rifacimento completo come quota del valore e $N_r$ il numero di anni fra due rifacimenti, cosicché il terzo addendo è l'accantonamento annuo. L'utile netto sottrae l'imposta sul canone:
\begin{equation}
U \;=\; \mathrm{NOI} \;-\; \Theta(L_e),
\label{eq:utile}
\end{equation}

\paragraph{In parole.} Dal reddito operativo netto si toglie l'imposta sul canone e resta l'utile netto, che e' il numero che finisce al numeratore del rendimento. La funzione $\Theta$ e' scritta con una lettera perche' cambia forma a seconda del regime fiscale scelto, e le sue quattro forme sono nella sezione seguente.

Sul caso del progetto, in cedolare secca al ventuno per cento: reddito operativo $1\,805$ euro, imposta $1\,120$, utile netto $684$ euro. Va notato che l'imposta e' calcolata sul canone e non sul reddito operativo: e' la ragione per cui su questo caso l'imposta si mangia quasi due terzi del reddito operativo, un rapporto che sorprende chi si aspetta di pagare il ventuno per cento su quello che gli resta.
con $\Theta$ funzione d'imposta che dipende dal regime scelto.

\begin{osservazione}[Un costo sta in un posto solo]
La \eqref{eq:noi} include l'accantonamento per la ristrutturazione, e la proiezione del flusso di cassa del capitolo seguente lo eredita attraverso $\mathrm{NOI}$ invece di ricalcolarlo. È la conseguenza di un difetto reale del modello, in cui la voce compariva sia nel conto economico sia come colonna autonoma della proiezione, e veniva quindi sottratta due volte: il flusso risultava peggiore del vero di quell'importo, ogni anno, e nessun controllo lo segnalava perché entrambe le formule erano corrette prese da sole.
\end{osservazione}

\subsection{Le quattro funzioni d'imposta}

I quattro regimi si distinguono per la forma di $\Theta$, e le differenze non sono solo di aliquota.
\begin{align}
\Theta_{\text{ced}}(L_e) &= L_e\,\theta_c, \label{eq:ced}\\
\Theta_{\text{conc}}(L_e) &= L_e\,\theta_{cc}, \label{eq:conc}\\
\Theta_{\text{ord}}(L_e) &= L_e\,(1-\alpha_b)\,(\theta_m + \theta_a) \;+\; \frac{\max(L_e\,\tau_\ell,\;\tau_{\ell,\min})}{2}, \label{eq:ord}\\
\Theta_{\text{breve}}(L_e) &= L_e\,\theta_b. \label{eq:breve}
\end{align}

\paragraph{In parole.} Le quattro righe sono i quattro modi di tassare lo stesso affitto, e le prime due sono semplici: si applica un'aliquota fissa al canone effettivo.

La terza, il regime ordinario, e' composta di due pezzi. Il primo dice che il canone entra nel reddito complessivo dopo un abbattimento forfettario, e viene tassato con l'aliquota marginale del contribuente piu' le addizionali locali: qui il conto dipende da chi sei, non dall'immobile. Il secondo pezzo e' l'imposta di registro annuale sul contratto, dovuta nel solo regime ordinario e divisa a meta' fra proprietario e inquilino, con un minimo di legge da cui viene il $\max$.

La quarta, la locazione breve, applica un'aliquota fissa come le prime due, ma il canone a cui si applica e' un ricavo di natura diversa, perche' viene da tariffe per notte e non da un contratto annuale.

Il confronto fra le quattro non e' quindi un confronto fra aliquote: cambia anche la base, cambia il fatto che l'imposta di registro annuale sparisca in tre casi su quattro, e in un caso su quattro il risultato dipende dalla posizione fiscale personale di chi compra.

Nella \eqref{eq:ord}, $\alpha_b$ è l'abbattimento forfettario, $\theta_m$ l'aliquota marginale IRPEF del contribuente, $\theta_a$ le addizionali locali, e l'ultimo addendo è l'imposta di registro annuale sul contratto, ripartita a metà fra locatore e conduttore. La cedolare secca la sostituisce, ed è una parte del suo vantaggio che l'aliquota nuda non mostra.

Due proprietà strutturali distinguono i regimi al di là dell'aliquota. La prima è che $\Theta_{\text{ord}}$ dipende da $\theta_m$, cioè dalla posizione fiscale personale: lo stesso immobile ha regimi ottimali diversi per due contribuenti diversi, e il confronto fra regimi non è una proprietà dell'immobile. La seconda è che la cedolare comporta la rinuncia all'aggiornamento ISTAT del canone, cioè $g=0$ nella notazione del capitolo sull'inflazione, e questa è una conseguenza sull'\emph{intera successione} dei canoni futuri e non sull'imposta di un anno. Il capitolo \ref{sec:inflazione} la quantifica.

\subsection{Proiezione del flusso di cassa}

Con indicizzazione $g$ del canone, la proiezione annuale è
\begin{equation}
F_t \;=\; \underbrace{L_e\,(1+g)^{t-1} - \big(\mathrm{NOI}\text{-costi}\big)(1+g)^{t-1} - \Theta_t}_{\text{gestione}} \;-\; \underbrace{\Gamma_{\text{annua}}\,\mathbb{1}[t\le n/12]}_{\text{debito}},
\qquad t=1,\dots,N,
\label{eq:proiezione}
\end{equation}
e al termine si aggiunge il flusso di uscita
\begin{equation}
F_N^{\text{uscita}} \;=\; P\,(1+\rho)^N (1-\xi) \;-\; D_{\min(N,\,n/12)} \;-\; \Pi,
\label{eq:uscita}
\end{equation}
con $\xi$ costi di vendita in quota e $\Pi$ eventuale imposta sulla plusvalenza. La \eqref{eq:uscita} è ciò che rende il tasso interno del capitolo seguente confrontabile con un investimento finanziario: senza l'uscita si confronterebbe una rendita con un capitale.

\section{Gli indicatori, e cosa ciascuno misura}

\subsection{Definizioni}

\begin{align}
\text{rendimento lordo} \quad & r_g = \frac{L}{P}, \label{eq:lordo}\\
\text{rendimento netto} \quad & r_n = \frac{U}{C}, \label{eq:netto}\\
\text{cap rate} \quad & r_c = \frac{\mathrm{NOI}}{C}, \label{eq:cap}\\
\text{cash on cash} \quad & r_e = \frac{U - \Gamma_{\text{annua}}}{E}, \label{eq:coc}\\
\text{DSCR} \quad & \mathrm{DSCR} = \frac{\mathrm{NOI}}{\Gamma_{\text{annua}}}. \label{eq:dscr}
\end{align}

\paragraph{In parole.} Cinque rapporti, che vanno letti guardando cosa hanno al numeratore e cosa al denominatore, perche' e' li' che sta tutta la differenza.

Il rendimento lordo mette il canone dichiarato sopra il prezzo. Non tiene conto di nessun costo e usa il prezzo invece del costo totale: e' il numero degli annunci, ed e' sistematicamente il piu' alto dei cinque.

Il rendimento netto mette l'utile, cioe' il canone dopo tutti i costi e le imposte, sopra il costo totale, cioe' il prezzo piu' tutto quello che si e' pagato per comprare. E' il numero da usare per decidere.

Il cap rate mette il reddito operativo, quindi prima delle imposte sul reddito e senza considerare il mutuo, sopra il costo totale. Serve a confrontare due immobili fra loro senza che il modo in cui li si finanzia sporchi il confronto.

Il cash on cash mette la cassa che resta dopo aver pagato anche la rata, sopra il capitale proprio effettivamente immobilizzato. Misura l'operazione e non l'immobile, e con la leva puo' essere negativo mentre gli altri quattro sono positivi.

Il debt service coverage ratio non e' un rendimento ma un rapporto di copertura: quante volte il reddito operativo copre la rata. Sotto uno, la differenza esce dalla tasca del proprietario ogni mese.

Sul caso del progetto i cinque valgono, nell'ordine: cinque per cento, $0{,}52$ per cento, $1{,}37$ per cento, meno $10{,}95$ per cento, e $0{,}34$. Sono cinque numeri della stessa operazione, e chi ne cita uno solo sceglie quale storia raccontare.

I cinque non sono varianti dello stesso numero: hanno numeratori e denominatori diversi perché rispondono a domande diverse. Il rendimento lordo è quello degli annunci ed è il meno informativo, perché ignora tutti i costi e usa $P$ al denominatore; serve a scremare. Il cap rate esclude il debito e le imposte sul reddito, quindi confronta immobili fra loro a prescindere da come sono finanziati. Il rendimento netto misura l'immobile per l'investitore che paga le imposte. Il \emph{cash on cash} misura l'operazione finanziata, e il suo denominatore è il capitale proprio. Il DSCR non è un rendimento ma un rapporto di copertura, e la sua soglia è uno.

\begin{osservazione}
Fra $r_g$ e $r_n$ si perdono di norma due punti e mezzo, e chi promette un netto vicino al lordo sta contando male o non sta contando le stesse cose. La verifica rapida è chiedere quale sia il denominatore e se lo sfitto sia nel numeratore.
\end{osservazione}

\subsection{L'effetto della leva}

\begin{proposizione}[Rendimento del capitale proprio]
Sia $r_c$ il cap rate, $c_d$ il costo del debito come quota del debito, $D=M$ il debito e $E=C-M$ il capitale proprio. Ignorando le imposte sul reddito, il rendimento del capitale proprio soddisfa
\begin{equation}
r_E \;=\; r_c \;+\; \frac{D}{E}\,\big(r_c - c_d\big).
\label{eq:leva}
\end{equation}

\paragraph{In parole.} La formula dice che il rendimento del proprio capitale e' il rendimento dell'immobile, piu' un premio che dipende da due cose: quanto debito si ha rispetto al capitale proprio, e quanto l'immobile rende in piu' rispetto a quanto costa il denaro.

Il pezzo fra parentesi e' la differenza fra i due tassi, e il fattore che lo moltiplica e' la leva. Se l'immobile rende piu' del debito, la parentesi e' positiva e la leva la amplifica: piu' debito, piu' rendimento sul proprio. Se l'immobile rende meno del debito, la parentesi e' negativa e la leva amplifica una perdita.

Sul caso del progetto: cap rate $1{,}37$ per cento, costo del debito $3{,}2$ per cento, quindi la parentesi vale meno $1{,}83$ punti. Con un debito di novantamila su un capitale proprio di $41\,557$, il rapporto $D/E$ vale $2{,}17$, e il rendimento del capitale proprio scende di quasi quattro punti sotto il cap rate. Non e' un difetto del calcolo: e' la leva che lavora al contrario, ed e' la ragione per cui in questo caso il cash on cash e' negativo mentre il cap rate e' positivo.
\end{proposizione}

\begin{proof}
Il reddito dell'attivo è $r_c\,C = r_c\,(D+E)$; il costo del debito è $c_d D$; ciò che resta al capitale proprio è la differenza, quindi
\[
r_E = \frac{r_c(D+E) - c_d D}{E} = r_c\,\frac{D+E}{E} - c_d\,\frac{D}{E}
= r_c + \frac{D}{E}\,(r_c - c_d),
\]
dove l'ultimo passaggio raccoglie $r_c$.
\end{proof}

La \eqref{eq:leva} è il risultato più utile del capitolo, perché rende esplicita una condizione che di solito si dà per ovvia e non lo è: la leva aumenta il rendimento del capitale proprio \emph{se e solo se} $r_c > c_d$, cioè se l'immobile rende più di quanto costa il denaro. Se $r_c < c_d$ la leva amplifica una perdita, e lo fa in proporzione a $D/E$. Sul caso di riferimento del progetto $r_c = 1{,}37$ per cento e il costo del debito è $3{,}2$ per cento: la disuguaglianza è rovesciata, e infatti il \emph{cash on cash} risulta negativo mentre il cap rate è positivo. Non è un difetto del calcolo, è la leva che lavora contro.

\begin{osservazione}[Perché il modello non usa la \eqref{eq:leva} per calcolare]
La \eqref{eq:leva} vale con $c_d$ costante e senza imposte, e in un mutuo ad ammortamento $c_d$ non è costante, perché la rata contiene anche capitale. Il modello calcola quindi $r_e$ direttamente dalla \eqref{eq:coc}, con il flusso di cassa effettivo, e usa la \eqref{eq:leva} come strumento di lettura: serve a capire il segno e l'ordine di grandezza, non a produrre il numero.
\end{osservazione}

\subsection{Tasso interno di rendimento}

\begin{definizione}
Il tasso interno di rendimento della successione $F_0,\dots,F_N$ è una radice dell'equazione
\begin{equation}
\VAN(y) \;=\; \sum_{t=0}^{N} \frac{F_t}{(1+y)^t} \;=\; 0 .
\label{eq:tir}
\end{equation}

\paragraph{In parole.} La formula non e' un calcolo ma una domanda: qual e' il tasso di sconto che rende nullo il valore attuale di questa operazione. Si legge cosi': trova il numero $y$ tale che, scontando tutti i flussi a quel tasso, la somma faccia esattamente zero.

Il significato e' che $y$ e' il rendimento che l'operazione produce, misurato in modo confrontabile con qualsiasi altro investimento. Se il tasso interno e' il quattro per cento, l'operazione rende come un titolo al quattro per cento che paghi gli stessi flussi negli stessi momenti.

Non esiste una formula chiusa per ricavare $y$: si cerca per tentativi, e il modo con cui il modello lo fa e' descritto poco sotto.
\end{definizione}

Tre questioni vanno affrontate, perché il TIR è l'indicatore su cui si sbaglia più spesso.

\paragraph{Esistenza e unicità.} La \eqref{eq:tir} è un polinomio in $v=(1+y)^{-1}$ di grado $N$, e in generale ha $N$ radici complesse. La regola dei segni di Cartesio garantisce che il numero di radici reali positive non superi il numero di cambi di segno della successione dei flussi. Per una struttura tipica di questo dominio, cioè un esborso iniziale negativo seguito da flussi di segno costante, il cambio di segno è unico e la radice reale positiva è unica: il TIR è ben definito. Quando la successione cambia segno più volte, per esempio con una ristrutturazione straordinaria a metà orizzonte, possono esistere più radici e il TIR perde significato. Il modello non pretende di rilevare questa condizione, e la dichiara come limite.

\paragraph{Metodo numerico.} Non esiste forma chiusa per $N>4$. Il motore usa la bisezione sull'intervallo $[-0{,}9999,\,10]$, con tolleranza $10^{-7}$ e duecento iterazioni. La bisezione converge linearmente, dimezzando l'ampiezza dell'intervallo a ogni passo, quindi duecento iterazioni sono largamente sufficienti; è preferita a Newton, che converge quadraticamente ma può divergere quando la derivata è piccola, situazione tipica proprio dei profili immobiliari, dove il primo termine è un esborso grande e i successivi sono piccoli e dello stesso segno. Se non esiste cambio di segno nell'intervallo, la funzione restituisce zero invece di un valore inventato.

\paragraph{Che cosa il TIR non dice.} Il TIR è un tasso di rendimento, non una misura di rischio, e assume implicitamente il reinvestimento dei flussi intermedi allo stesso tasso. Due operazioni con lo stesso TIR e scale diverse non sono equivalenti, ed è la ragione per cui il modello riporta accanto il valore attuale netto \eqref{eq:van}, che è additivo e commensurabile, e il multiplo sul capitale proprio.

\subsection{Relazione fra VAN e TIR}

Il VAN è funzione decrescente del tasso di sconto quando i flussi hanno un solo cambio di segno, e il TIR è per definizione il tasso che lo annulla. Quindi $\VAN(s)>0 \iff s< \TIR$ nel caso regolare: le due regole di decisione coincidono. Non coincidono nel confronto fra progetti di scala diversa, dove il VAN va usato e il TIR no, perché il primo è additivo e il secondo è un tasso.

\section{L'inversione: il prezzo massimo sostenibile}
\label{sec:inversione}

\subsection{Il problema}

La domanda della trattativa non è quanto rende l'immobile a un prezzo dato, ma quale prezzo, al massimo, giustifica l'operazione a un rendimento netto dichiarato accettabile $r^\dagger$. Formalmente si cerca $P^\star$ tale che
\begin{equation}
\frac{U(P^\star)}{C(P^\star)} \;=\; r^\dagger .
\label{eq:problema-inversione}
\end{equation}

\paragraph{In parole.} La domanda scritta in simboli e': trova il prezzo $P^\star$ per cui il rendimento netto e' esattamente pari all'obiettivo che ti sei dato. E' la domanda della trattativa, e non e' la stessa cosa che chiedersi quanto rende l'immobile al prezzo richiesto: qui il prezzo e' l'incognita e il rendimento e' il dato.

Il simbolo con la stella indica un valore particolare della grandezza, quello che risolve il problema, e serve a distinguerlo dal prezzo generico $P$ che compare nelle due funzioni.

Il modo sbagliato di risolverla, e quello che il modello adottava fino alla revisione del 1 settembre 2026, è dividere il costo sostenibile per uno più l'incidenza dei costi accessori misurata allo scenario base:
\begin{equation}
P^\star_{\text{errato}} \;=\; \frac{U/r^\dagger}{1 + A(P_0)/P_0}.
\label{eq:inversione-errata}
\end{equation}

La \eqref{eq:inversione-errata} sbaglia per due ragioni indipendenti che si sommano nello stesso verso.

\paragraph{Prima causa: l'incidenza non è costante.} Dalla \eqref{eq:costo-affine} si ha $A(P)/P = k + c/P$, che è una funzione decrescente di $P$ e non un parametro. Congelarla al valore dello scenario base significa usare $k + c/P_0$ dove servirebbe $k + c/P^\star$, e poiché $P^\star < P_0$ nel caso interessante, si sottostima l'incidenza vera.

\paragraph{Seconda causa: l'utile non è indipendente dal prezzo.} Dalla \eqref{eq:noi} si vede che manutenzione e accantonamento per la ristrutturazione sono quote del valore:
\begin{equation}
U(P) \;=\; U_0 \;-\; (P-P_0)\,\mathrm{m}^\dagger,
\qquad
\mathrm{m}^\dagger := \mathrm{m} + \frac{\mathrm{r}}{N_r},
\label{eq:utile-di-p}
\end{equation}

\paragraph{In parole.} L'utile netto non e' un numero fisso: cambia se cambia il prezzo. La formula dice di quanto: partendo dall'utile calcolato al prezzo attuale, per ogni euro di prezzo in piu' l'utile scende di $m^\dagger$ euro.

La ragione e' che due voci di costo sono calcolate come quota del valore dell'immobile, la manutenzione ordinaria e l'accantonamento per la ristrutturazione. Se l'immobile costa piu' caro, manutenerlo e rifarlo costa piu' caro. Sul caso del progetto $m^\dagger$ vale $1{,}825$ per cento, cioe' un per cento di manutenzione piu' $0{,}825$ di accantonamento: comprando ventimila euro in meno, l'utile annuo sale di trecentosessantacinque euro.
dove $U_0$ è l'utile al prezzo base. Quindi abbassare il prezzo \emph{alza} l'utile, e la \eqref{eq:inversione-errata} lo tiene fermo.

\subsection{La soluzione esatta}

\begin{proposizione}[Prezzo massimo sostenibile]
Sul tratto in cui $C$ è affine, cioè fuori dalla regione in cui il minimo di legge dell'imposta di registro è vincolante, la \eqref{eq:problema-inversione} ha l'unica soluzione
\begin{equation}
\boxed{\;
P^\star \;=\; \frac{U_0 + P_0\,\mathrm{m}^\dagger - r^\dagger c}{r^\dagger (1+k) + \mathrm{m}^\dagger}\;}
\label{eq:pmax}
\end{equation}

\paragraph{In parole.} La formula e' il risultato dell'inversione, e si legge per pezzi.

Al numeratore ci sono tre termini. Il primo e' l'utile al prezzo attuale. Il secondo aggiunge il risparmio sui costi proporzionali che si otterrebbe comprando a un prezzo qualunque, calcolato sul prezzo attuale: e' un termine tecnico che nasce dallo spostare l'incognita a destra. Il terzo sottrae il rendimento obiettivo applicato ai costi fissi, cioe' la parte di rendimento che i costi fissi si mangiano comunque.

Al denominatore c'e' il rendimento obiettivo moltiplicato per il fattore di costo totale, piu' i costi annui che scalano col prezzo. E' quanto rendimento obiettivo costa ogni euro di prezzo.

Sul caso del progetto: numeratore $684 + 2\,190 - 287 = 2\,587$, denominatore $0{,}0415 + 0{,}0183 = 0{,}0597$, risultato $43\,445$ euro. Cioe': un immobile che rende $684$ euro netti l'anno giustifica un prezzo di circa quarantatremila euro se si pretende un rendimento netto del quattro per cento. Al prezzo trattato di centoventimila il rendimento e' lo $0{,}52$ per cento, e la distanza fra i due prezzi e' lo sconto che servirebbe.
\end{proposizione}

\begin{proof}
Si sostituiscono la \eqref{eq:costo-affine} e la \eqref{eq:utile-di-p} nella \eqref{eq:problema-inversione}:
\[
U_0 - (P-P_0)\,\mathrm{m}^\dagger \;=\; r^\dagger\big(P(1+k)+c\big).
\]
Si espande e si raccolgono i termini in $P$ a destra:
\[
U_0 + P_0 \mathrm{m}^\dagger \;=\; r^\dagger(1+k)P + r^\dagger c + P\,\mathrm{m}^\dagger
\;=\; P\big(r^\dagger(1+k)+\mathrm{m}^\dagger\big) + r^\dagger c,
\]
da cui la tesi dividendo per il coefficiente di $P$, che è positivo per $r^\dagger>0$ e $k,\mathrm{m}^\dagger \ge 0$, il che garantisce anche l'unicità.
\end{proof}

L'equazione è di primo grado, quindi la soluzione è chiusa e non richiede metodi numerici. I tre coefficienti $k$, $c$ e $\mathrm{m}^\dagger$ stanno nel foglio in tre celle visibili con la loro nota, e non dentro la formula finale, perché un risultato di cui non si possono riconoscere i pezzi non si può contestare: sul caso di riferimento valgono $3{,}7$ per cento, $7\,165$ euro e $1{,}8$ per cento, e chi conosce il dominio riconosce nel primo la provvigione con la sua IVA, nel secondo il fatto che il prezzo-valore è attivo perché l'intera imposta di registro è finita nella parte fissa, e nel terzo l'uno per cento di manutenzione più lo zero virgola ottantadue di accantonamento.

\subsection{Quanto valeva l'errore}

Sul caso di riferimento, con $U_0 = 684$ euro, $r^\dagger = 4$ per cento, $P_0 = 120\,000$ euro:
\begin{align*}
P^\star_{\text{errato}} &= \frac{684/0{,}04}{1{,}0963} = \frac{17\,112}{1{,}0963} = 15\,609 \text{ euro},\\
P^\star &= \frac{684 + 120\,000\cdot 0{,}01825 - 0{,}04\cdot 7\,165}{0{,}04\cdot 1{,}0366 + 0{,}01825} = \frac{2\,587{,}4}{0{,}059714} = 43\,445 \text{ euro}.
\end{align*}
Un fattore prossimo a tre, nel verso che fa apparire impossibile qualunque trattativa. Per un numero il cui unico scopo è dire quanto offrire, è il modo peggiore di sbagliare.

\subsection{Il controllo di chiusura}

La \eqref{eq:pmax} è esatta solo sul tratto affine. Il modello affianca al risultato una cella che ricalcola il rendimento a $P^\star$ con le formule esatte, minimo di legge compreso, e mostra lo scarto dalla soglia:
\begin{equation}
\varepsilon \;=\; \frac{U_0 - (P^\star-P_0)\mathrm{m}^\dagger}{P^\star + T_a(P^\star) + P^\star p(1+\upsilon_p) + \nu + \omega + \Omega_M} \;-\; r^\dagger .
\label{eq:controllo}
\end{equation}
Sul caso di riferimento $\varepsilon = 0$ a quattro decimali. Il controllo sta nel foglio e non in un test per una ragione precisa: il caso che rompe la linearità si presenta quando chi usa il file cambia rendita, aliquota o obiettivo a video, non nel caso precaricato che un test coprirebbe.

\begin{osservazione}[La lezione generalizzabile]
Un'incidenza percentuale misurata su uno scenario non è un parametro del modello: è un rapporto fra due grandezze che dipendono entrambe dalla variabile che si sta muovendo. Usarla come costante è il modo più comune di introdurre un errore che non si vede, perché il risultato resta dell'ordine di grandezza giusto. La forma corretta di un'inversione è scomporre in parte proporzionale e parte fissa e risolvere l'equazione, non stimare un rapporto.
\end{osservazione}

\section{Sensibilità e rischio}

\subsection{Derivate e elasticità}

La sensibilità del rendimento netto a una variabile $\psi$ è la derivata parziale $\partial r_n/\partial \psi$, e la sua forma adimensionale è l'elasticità
\begin{equation}
\epsilon_\psi \;=\; \frac{\partial r_n}{\partial \psi}\cdot\frac{\psi}{r_n},
\label{eq:elasticita}
\end{equation}
cioè la variazione percentuale del rendimento per una variazione dell'uno per cento della variabile. Le elasticità sono la grandezza giusta per confrontare l'importanza di variabili con unità diverse, che è esattamente il problema del diagramma a tornado.

Sul canone si ha, tenendo tutto il resto fermo,
\begin{equation}
\frac{\partial U}{\partial L} = \Big(1-\frac{\sigma}{12}\Big)(1-\beta)\,(1-\theta),
\qquad
\frac{\partial r_n}{\partial L} = \frac{1}{C}\,\frac{\partial U}{\partial L},
\label{eq:sens-canone}
\end{equation}
che mostra una cosa non ovvia: l'effetto di un aumento del canone sul rendimento è attenuato tre volte, dallo sfitto, dalla morosità e dall'imposta. Con $\sigma=1$, $\beta=0{,}03$ e $\theta=0{,}21$ il fattore complessivo è $0{,}917\cdot 0{,}97\cdot 0{,}79 = 0{,}703$: settanta centesimi di ogni euro di canone in più arrivano al rendimento.

Sul prezzo la derivata è più ricca, perché il prezzo entra due volte:
\begin{equation}
\frac{\partial r_n}{\partial P} \;=\; \frac{-\mathrm{m}^\dagger\,C - U\,(1+k)}{C^2} \;<\; 0 ,
\label{eq:sens-prezzo}
\end{equation}
dove il primo addendo del numeratore viene dai costi che scalano col prezzo e il secondo dal costo totale al denominatore. La \eqref{eq:sens-prezzo} è sempre negativa: nessuna sorpresa, ma la sua scomposizione dice che circa un terzo dell'effetto passa dal numeratore, cosa che una tabella di sensibilità non mostra.

\subsection{Il diagramma a tornado}

Il tornado ordina le variabili per ampiezza dell'effetto su una grandezza scelta, applicando a ciascuna uno scostamento e tenendo le altre al valore base. Se $\psi_h$ è la variabile $h$ e $\Delta_h$ lo scostamento applicato, la barra è
\begin{equation}
b_h \;=\; \big|\,G(\psi_h + \Delta_h) - G(\psi_h - \Delta_h)\,\big|,
\label{eq:tornado}
\end{equation}
con $G$ la grandezza osservata. Due avvertenze, entrambe dichiarate nel foglio. Gli scostamenti $\Delta_h$ devono essere commensurabili fra loro, altrimenti l'ordinamento misura la scelta degli scostamenti e non l'importanza delle variabili; e il tornado è per costruzione univariato, quindi non vede le interazioni: due variabili che singolarmente contano poco possono contare molto insieme.

\subsection{Simulazione: perché non basta il caso base}

Tre scenari scelti a mano rispondono alla domanda sbagliata, perché non dicono quanto sia probabile ciascuno. La simulazione sostituisce gli scenari con una distribuzione di esiti.

\paragraph{Impianto.} Si estraggono $S=1000$ realizzazioni di un vettore di variabili normali standard $Z = (Z_1,\dots,Z_5)$ e una uniforme $V\sim \mathcal{U}(0,1)$ per gli eventi discreti. Ogni grandezza incerta si costruisce come trasformazione della corrispondente normale:
\begin{align}
L^{(s)} &= \max\big(0,\; L\,(1+Z_1^{(s)}\,\varsigma_L)\big),\\
\sigma^{(s)} &= \mathrm{clip}\big(\sigma + Z_2^{(s)}\,\varsigma_\sigma,\; 0,\; 12\big),\\
i^{(s)} &= \max\big(0,\; i + Z_3^{(s)}\,\varsigma_i\big),\\
\rho^{(s)} &= \rho + Z_4^{(s)}\,\frac{\varsigma_\rho}{\sqrt{N}} .
\label{eq:estrazioni}
\end{align}

\paragraph{La correzione che vale più della tecnica.} L'ultima riga della \eqref{eq:estrazioni} contiene una divisione per $\sqrt{N}$ che non è cosmetica. La prima versione del modello trattava l'estrazione sulla rivalutazione come un regime permanente per tutto l'orizzonte, cioè applicava $\rho + Z\varsigma_\rho$ per $N$ anni consecutivi. È sbagliato: la rivalutazione si compone, quindi la grandezza incerta rilevante non è la variazione di un anno ma la \emph{media} del periodo, e la deviazione standard della media di $N$ variabili indipendenti con deviazione $\varsigma$ vale $\varsigma/\sqrt{N}$. Senza la correzione la coda alta produceva patrimoni finali fuori scala: una simulazione matematicamente coerente e finanziariamente assurda.

\paragraph{Riproducibilità.} Le estrazioni sono generate una volta sola alla costruzione del file, con un seme dichiarato, e scritte come numeri fissi in un foglio nascosto; il calcolo che sta sopra è invece formula viva. La ragione è una proprietà della funzione casuale nativa dei fogli di calcolo: è volatile, quindi ogni ricalcolo rigenera tutti i numeri, due letture consecutive dello stesso file danno risultati diversi e nulla è verificabile. Separando estrazione e calcolo si ottiene una simulazione insieme stabile e interattiva, che nessuna delle due forme pure dava.

\paragraph{Il limite dichiarato.} Le variabili sono estratte indipendenti, mentre nella realtà tassi, prezzi, sfitto e morosità si muovono insieme. Introdurre una struttura di correlazione richiederebbe di stimare una matrice che nessuno ha, e sostituirebbe un'assunzione dichiarata con una nascosta. Il modello resta indipendente e lo dice.

\paragraph{Errore di campionamento.} Con $S$ estrazioni, la stima di una probabilità $p$ ha errore standard $\sqrt{p(1-p)/S}$, che per $p\approx 0{,}5$ e $S=1000$ vale circa $1{,}6$ punti percentuali. Le probabilità riportate dal modello vanno quindi lette a due cifre e non a tre, e questo è il motivo per cui il foglio non stampa decimali su di esse.

\section{Inflazione applicata a questa operazione}
\label{sec:inflazione}

\subsection{Perché il segno non è ovvio}

L'inflazione non agisce nello stesso verso su tutte le componenti, e in un acquisto a leva i versi si compensano solo in parte. Vale scomporre.

\begin{table}[H]
\centering
\small
\begin{tabular}{@{}lll@{}}
\toprule
Componente & Natura & Effetto dell'inflazione \\
\midrule
Debito residuo & importo nominale fisso & favorevole: si erode \\
Rata a tasso fisso & importo nominale fisso & favorevole: si alleggerisce \\
Rata a tasso variabile & indicizzata all'indice & sfavorevole: cresce \\
Canone, cedolare secca & non indicizzabile, $g=0$ & sfavorevole: perde tutto $\pi$ \\
Canone, regime ordinario & indicizzabile ISTAT & neutro se $g=\pi$ \\
Valore dell'immobile & nominale, cresce di $\rho$ & neutro se $\rho=\pi$ \\
Costi operativi & nominali, crescono & sfavorevole se non si ribaltano \\
\bottomrule
\end{tabular}
\caption{Direzione dell'effetto dell'inflazione per componente.}
\end{table}

\subsection{Il guadagno del debitore}

\begin{proposizione}
Il valore, in potere d'acquisto di oggi, dell'erosione operata dall'inflazione su un debito residuo nominale $D_N$ a un orizzonte di $N$ anni, è
\begin{equation}
\Delta_\pi \;=\; D_N\Big(1 - (1+\pi)^{-N}\Big).
\label{eq:regalo}
\end{equation}

\paragraph{In parole.} Il debito e' un impegno a consegnare un certo numero di euro, non un certo potere d'acquisto. Se fra dieci anni si devono ancora quarantamila euro, quei quarantamila euro compreranno meno cose di quante ne comprerebbero oggi, e la differenza e' un guadagno per chi deve pagarli.

La formula lo quantifica: prendi il debito residuo, e sottraine il suo valore in euro di oggi. Con quarantamila euro di residuo fra venticinque anni e un'inflazione del due per cento, il valore attuale e' $40\,000/1{,}02^{25} = 24\,380$ euro, quindi l'inflazione ha eroso circa $15\,600$ euro di debito reale.

Vale solo a tasso fisso, e la ragione e' semplice: su un variabile l'inflazione fa salire il tasso di riferimento, quindi si paga di piu' in interessi e il guadagno sul capitale si compensa in tutto o in parte con la perdita sul servizio del debito.
\end{proposizione}

\begin{proof}
Il debito è un'obbligazione a consegnare $D_N$ euro nominali. Il suo valore in euro di oggi è $D_N(1+\pi)^{-N}$ per la \eqref{eq:deflaziona}. La differenza fra l'importo nominale dovuto e il suo valore reale è la tesi.
\end{proof}

La \eqref{eq:regalo} è un trasferimento reale di ricchezza dal creditore al debitore, e in un mutuo residenziale il debitore è chi compra. È la ragione per cui l'inflazione, in un acquisto a leva a tasso fisso, ha un lato favorevole che il rendimento nominale non mostra e quello reale mostra solo in aggregato.

\begin{osservazione}[Vale solo a tasso fisso]
Su un variabile l'inflazione fa salire l'indice di riferimento, la rata cresce e lo sconto sul debito si paga in interessi. La \eqref{eq:regalo} continua a valere aritmeticamente ma è compensata, in tutto o in parte, dall'aumento del servizio del debito: il capitolo sul percorso del tasso a gradini serve esattamente a quantificare il secondo termine.
\end{osservazione}

\subsection{Rendimenti reali}

Applicando la \eqref{eq:fisher} agli indicatori si ottiene
\begin{equation}
r_n^* = \frac{1+r_n}{1+\pi}-1,
\qquad
\TIR^* = \frac{1+\TIR}{1+\pi}-1,
\qquad
\rho^* = \frac{1+\rho}{1+\pi}-1 .
\label{eq:reali}
\end{equation}

La terza è quella che merita attenzione: se si assume $\rho = \pi$, cioè che l'immobile si rivaluti quanto l'inflazione, allora $\rho^*=0$ esattamente, e il valore finale in potere d'acquisto coincide con il prezzo pagato. Non è un'assunzione pessimistica: in termini reali il mercato residenziale italiano è rimasto sostanzialmente fermo per un ventennio, e porre $\rho=\pi$ è già benevolo.

La variazione reale del canone è
\begin{equation}
g^* = \frac{1+g}{1+\pi}-1,
\label{eq:canone-reale}
\end{equation}
che con la cedolare secca, dove $g=0$ per obbligo, vale $-\pi/(1+\pi)$: il canone perde in termini reali quasi tutta l'inflazione, ogni anno, per tutta la durata dell'opzione.

\subsection{Il costo dell'indicizzazione rinunciata}

Questa è la parte del capitolo che quantifica una scelta fiscale che di solito si fa guardando solo l'aliquota.

\begin{proposizione}
Il valore attuale del canone a cui si rinuncia non indicizzando, su un orizzonte di $N$ anni, al tasso di sconto $s$, vale
\begin{equation}
\Lambda \;=\; L\,\Big(F(\pi,s,N) - F(g,s,N)\Big),
\label{eq:rinuncia}
\end{equation}

\paragraph{In parole.} La formula calcola quanto costa, in euro di oggi, rinunciare ad aggiornare il canone. Si legge: prendi il canone iniziale, moltiplicalo per la differenza fra il fattore di una rendita che cresce all'inflazione e il fattore di una rendita che cresce come la tua.

Se non si indicizza affatto, cioe' $g = 0$, la seconda rendita e' un canone fermo e la differenza fra i due fattori e' tutto il valore dell'indicizzazione. Sul caso del progetto, con un canone di seimila euro l'anno, inflazione al due per cento e sconto al tre su venticinque anni: $6\,000 \times (21{,}64 - 17{,}41) = 25\,384$ euro. Al netto dell'imposta sul canone scendono a circa ventimila.

E' un costo che non si vede mai, perche' non e' una spesa: e' un canone che non e' stato chiesto.
con $F$ il fattore \eqref{eq:crescente}. Al netto dell'imposta sul canone la grandezza comparabile è $\Lambda(1-\theta)$.
\end{proposizione}

\begin{proof}
La successione dei canoni pienamente indicizzati è $L(1+\pi)^{t-1}$, quella effettiva $L(1+g)^{t-1}$. Il valore attuale della differenza è, per linearità,
\[
\sum_{t=1}^{N}\frac{L(1+\pi)^{t-1} - L(1+g)^{t-1}}{(1+s)^t}
= L\Big(\sum_{t=1}^{N}\frac{(1+\pi)^{t-1}}{(1+s)^t} - \sum_{t=1}^{N}\frac{(1+g)^{t-1}}{(1+s)^t}\Big),
\]
cioè la tesi per definizione di $F$.
\end{proof}

Il lato opposto della scelta è il risparmio d'imposta, il cui valore attuale su un canone effettivo costante è
\begin{equation}
\Sigma \;=\; \big(\Theta_{\text{ord}}(L_e) - \Theta_{\text{ced}}(L_e)\big)\,a_{N,s},
\label{eq:risparmio}
\end{equation}
e il saldo della scelta è $\Sigma - \Lambda(1-\theta)$.

\begin{osservazione}[Il risultato sul caso di riferimento, con le sue cautele]
Con $L=6\,000$ euro, $\pi=2$ per cento, $g=0$, $s=3$ per cento, $N=25$ anni si ottiene $F(\pi,s,N)=21{,}6437$ e $F(g,s,N)=17{,}4131$, quindi $\Lambda = 25\,384$ euro lordi e $20\,053$ euro al netto della cedolare al ventuno per cento. Il risparmio d'imposta attualizzato vale $12\,777$ euro. Il saldo è quindi \emph{negativo} per circa settemiladuecento euro: su quell'orizzonte la cedolare secca costa più di quanto risparmi.

Tre cautele, perché il confronto è fra grandezze non simmetriche e leggerlo come un verdetto sarebbe scorretto. L'orizzonte usato è quello del modello, decenni, mentre l'opzione per la cedolare si esercita per contratto e si può riconsiderare a ogni rinnovo: la cifra è il caso peggiore, quello di chi la rinnova sempre senza ripensarci. Il risparmio d'imposta dipende dall'aliquota marginale personale, che è un input e non un dato. E il canone concordato ha regole di aggiornamento proprie, fissate dall'accordo territoriale, che la \eqref{eq:rinuncia} non modella.

Ciò che il saldo dice con certezza è che la scelta ha due lati e che il secondo non è trascurabile. È un confronto che quasi nessuno fa, perché il risparmio si vede subito nella dichiarazione e la mancata indicizzazione non si vede mai, essendo un canone che non è stato chiesto.
\end{osservazione}

\section{Comprare oppure restare in affitto}

\subsection{Formalizzazione a parità di esborso}

Il confronto corretto non è fra la rata e il canone, che è il confronto che si fa spontaneamente e che ignora la componente di capitale della rata. Il confronto è fra due strategie che partono dallo stesso esborso iniziale.

Chi compra versa $E$ al tempo zero e poi, ogni anno, la rata più i costi del proprietario meno l'eventuale detrazione. Chi affitta investe $E$ in un portafoglio che rende $r_p$, paga il canone di mercato $L_a$ indicizzato all'inflazione, e investe ogni anno la differenza fra le uscite dell'altro e le proprie. Le due posizioni finali sono
\begin{align}
W_{\text{compra}} &= P(1+\rho)^N(1-\xi) - D_{\min(N,n/12)},
\label{eq:w-compra}\\
W_{\text{affitta}} &= \Big[E(1+r_p)^N + \sum_{t=1}^{N}\Delta_t (1+r_p)^{N-t}\Big] - \theta_p\,\max\big(0,\,\cdot\, - \text{versato}\big),
\label{eq:w-affitta}
\end{align}
con $\Delta_t$ la differenza fra le uscite annuali dei due e $\theta_p$ l'imposta sulle rendite finanziarie applicata alla sola plusvalenza.

\subsection{Le tre assunzioni che ribaltano l'esito}

L'esito dipende quasi solo da $r_p$, $\rho$ e $L_a$. La derivata rispetto a $r_p$ è la più violenta: cambiando $r_p$ di un punto percentuale l'esito spesso si rovescia. Questo dice quanto poco il verdetto vada preso come responso e quanto vada preso come mappa di sensibilità.

Due effetti non compaiono nelle \eqref{eq:w-compra} e \eqref{eq:w-affitta} e vanno tenuti a mente. Il primo è che la \eqref{eq:w-affitta} assume disciplina perfetta di chi affitta, cioè che investa davvero ogni euro di differenza, ogni anno, senza toccarlo; nella realtà quasi nessuno lo fa, e il mutuo funziona come piano di accumulo forzato. È un vantaggio comportamentale reale che il modello non sa misurare. Il secondo è che restano fuori la sicurezza abitativa, la libertà di intervenire sull'immobile, il rischio di sfratto e il vincolo di mobilità lavorativa: decisivi nella scelta di dove vivere, irrilevanti nella scelta di dove investire, ed è la ragione per cui le due domande vanno tenute separate anche quando riguardano lo stesso immobile.

\subsection{Il caso senza mutuo}

Se $M=0$ allora $E=C$ e le uscite di chi compra si riducono ai soli costi del proprietario. Il confronto continua a calcolare ma risponde a un'altra domanda: non più comprare a debito contro affittare e investire, ma immobilizzare il capitale in un immobile contro investirlo sui mercati. Per quella domanda l'indicatore corretto è il TIR reale confrontato con il rendimento reale del portafoglio, e il modello lo dichiara con un'avvertenza che compare da sola quando l'importo del mutuo è nullo.

\section{Il confronto fra immobili e la graduatoria}

\subsection{Normalizzazione}

Fra immobili di taglia diversa il prezzo non è confrontabile. La normalizzazione elementare è il prezzo per unità di superficie,
\begin{equation}
u \;=\; \frac{P}{S},
\label{eq:prezzo-mq}
\end{equation}
con $S$ superficie commerciale. La scelta della superficie commerciale, e non della calpestabile, non è indifferente: le quotazioni dell'Osservatorio del mercato immobiliare sono espresse su superficie commerciale, quindi il confronto con esse richiede la stessa base. Usare la calpestabile al numeratore e confrontare con quotazioni commerciali produce uno scarto sistematico verso l'alto del quindici-venti per cento.

\subsection{Lo scarto sulla quotazione di zona}

\begin{definizione}
Sia $[\underline{q},\overline{q}]$ l'intervallo di quotazione della zona omogenea per la tipologia dell'immobile. Lo scarto è
\begin{equation}
\delta \;=\; \frac{u}{\tfrac{1}{2}(\underline{q}+\overline{q})} - 1 .
\label{eq:scarto}
\end{equation}

\paragraph{In parole.} Prendi il prezzo al metro quadro dell'immobile, dividilo per il punto medio dell'intervallo di quotazione della sua zona, e togli uno. Il risultato dice di quanto il prezzo sta sopra o sotto la media della zona, in percentuale.

Con un prezzo al metro quadro di $2\,565$ euro e una zona quotata fra $1\,800$ e $3\,300$: il punto medio e' $2\,550$, il rapporto e' $1{,}006$, lo scarto e' piu' $0{,}6$ per cento. L'immobile e' praticamente in linea con la sua zona.

Tre avvertenze. Il punto medio e' una scelta e non un dato: la quotazione e' un intervallo proprio perche' nella zona ci sono immobili diversi, e schiacciarlo su un numero perde quell'informazione. Lo scarto non e' un giudizio, perche' la quotazione non vede stato, piano, affaccio ne' classe energetica. E se la zona non e' nota, il modello usa l'intervallo dell'intero Comune, che e' largo abbastanza da rendere lo scarto quasi privo di significato.
\end{definizione}

Tre proprietà vanno dichiarate. La \eqref{eq:scarto} usa la media dell'intervallo, che è una scelta e non un dato: la quotazione OMI è un intervallo perché la zona contiene immobili diversi, e collassarlo sulla media perde l'informazione sulla dispersione. Lo scarto è definito solo se la zona è nota: quando manca, il modello usa l'intervallo dell'intero Comune, che è largo e rende $\delta$ poco informativo, e lo segnala. E $\delta$ non è un giudizio: la quotazione è una media di zona per tipologia, quindi non vede stato di conservazione, piano, affaccio, classe energetica né lavori deliberati in condominio. Uno scarto molto negativo ha sempre una ragione, e la ragione va trovata prima di innamorarsene.

\subsection{Il prezzo che entra nel confronto}

Il foglio di confronto usa, per ogni riga, il prezzo obiettivo se compilato e il richiesto altrimenti:
\begin{equation}
P_{\text{riga}} \;=\; \begin{cases} P_{\text{obiettivo}} & \text{se } P_{\text{obiettivo}}>0,\\ P_{\text{richiesto}} & \text{altrimenti.}\end{cases}
\label{eq:prezzo-riga}
\end{equation}
Ne segue che lo scarto calcolato sul foglio di confronto differisce da quello calcolato sul registro, che usa sempre il richiesto. La differenza fra i due, letta al contrario, è lo sconto che si sta chiedendo: non è un'incoerenza ma un'informazione, purché si sappia quale delle due si sta leggendo, ed è la ragione per cui il modello calcola lo scarto due volte invece di leggerlo una volta dall'altro foglio.

\subsection{Il regime di acquisto è un dato della riga}

Le funzioni d'imposta \eqref{eq:imposte} dipendono da $\alpha$, agevolazione prima casa, e dal tipo di venditore. Nessuno dei due è una costante del modello: $\alpha$ non è nemmeno una caratteristica dell'immobile, è una caratteristica della posizione di chi compra rispetto a quell'immobile, e cambia tipicamente fra un immobile nel Comune di residenza e uno fuori. Applicare lo stesso regime a tutte le righe produce, come mostra l'osservazione del capitolo sul costo, una graduatoria ordinata al contrario nella parte alta. Il modello li dichiara per riga con tre stati, dove il terzo, il vuoto, significa eredita dal caso base: la neutralità sui registri già compilati non è cortesia verso il passato ma una proprietà verificabile.

\subsection{Perché non esiste un punteggio sintetico}

Sarebbe tecnicamente facile combinare gli indicatori in un unico numero, con pesi. Il modello non lo fa, e la ragione è che i pesi sarebbero arbitrari e la loro arbitrarietà scomparirebbe dentro il numero: un punteggio sintetico non aggiunge informazione, nasconde il ragionamento che l'ha prodotto. La graduatoria ordina per un criterio dichiarato, lo scarto sulla zona, e affianca gli altri indicatori senza fonderli; l'ordine in cui guardare gli immobili resta un campo compilato a mano, perché dipende anche da cose che il modello non sa, come la vicinanza a chi ci deve abitare o il fatto che il venditore abbia fretta.

\appendix

\section{Tavola di corrispondenza}

Ogni simbolo di questo documento corrisponde a un nome definito nel workbook e, dove esiste, a una funzione del motore Python. La tavola serve a rendere ogni numero rintracciabile fino alla sua formula.

\begin{table}[H]
\centering
\small
\begin{tabular}{@{}llll@{}}
\toprule
Simbolo & Formula & Nome nel workbook & Motore Python \\
\midrule
$V_c$ & \eqref{eq:catastale} & \texttt{valore\_catastale} & \texttt{valore\_catastale} \\
$B(P)$ & \eqref{eq:base} & \texttt{base\_registro} & \texttt{base\_imponibile\_registro} \\
$T_a$ & \eqref{eq:imposte} & \texttt{imposte\_totali} & \texttt{imposte\_acquisto} \\
$A$ & \eqref{eq:accessori} & \texttt{costi\_accessori} & \texttt{costo\_operazione} \\
$C$ & \eqref{eq:costo-affine} & \texttt{costo\_totale} & \texttt{costo\_operazione} \\
$E$ & --- & \texttt{esborso} & \texttt{costo\_operazione} \\
$\Gamma$ & \eqref{eq:rata} & \texttt{rata\_mensile} & \texttt{rata\_francese} \\
$D_k$ & \eqref{eq:residuo} & \texttt{debito\_residuo\_anno} & \texttt{piano\_ammortamento} \\
$L_e$ & \eqref{eq:ricavo-effettivo} & \texttt{ricavo\_effettivo} & \texttt{conto\_economico} \\
$\mathrm{NOI}$ & \eqref{eq:noi} & \texttt{noi\_annuo} & \texttt{conto\_economico} \\
$U$ & \eqref{eq:utile} & \texttt{utile\_locazione} & \texttt{conto\_economico} \\
$\Theta$ & \eqref{eq:ced}--\eqref{eq:breve} & --- & \texttt{imposta\_sul\_canone} \\
$r_n$ & \eqref{eq:netto} & --- & \texttt{metriche} \\
$r_c$ & \eqref{eq:cap} & --- & \texttt{metriche} \\
$r_e$ & \eqref{eq:coc} & --- & \texttt{metriche} \\
$\mathrm{DSCR}$ & \eqref{eq:dscr} & --- & \texttt{metriche} \\
$\TIR$ & \eqref{eq:tir} & \texttt{flussi\_tir} & \texttt{tir} \\
$\VAN$ & \eqref{eq:van} & \texttt{tasso\_sconto} & \texttt{van} \\
$P^\star$ & \eqref{eq:pmax} & --- & --- \\
$k$, $c$ & \eqref{eq:k}, \eqref{eq:c} & --- & --- \\
$F(g,s,n)$ & \eqref{eq:crescente} & --- & \texttt{fattore\_rendita\_crescente} \\
$r^*$ & \eqref{eq:fisher} & --- & \texttt{tasso\_reale} \\
$Y^*$ & \eqref{eq:deflaziona} & --- & \texttt{deflaziona} \\
$\Delta_\pi$ & \eqref{eq:regalo} & --- & \texttt{effetto\_inflazione} \\
$\Lambda$ & \eqref{eq:rinuncia} & --- & \texttt{effetto\_inflazione} \\
$\delta$ & \eqref{eq:scarto} & --- & \texttt{Annuncio.scarto\_su\_omi} \\
$u$ & \eqref{eq:prezzo-mq} & --- & \texttt{Annuncio.prezzo\_mq} \\
$i_k$ & \eqref{eq:gradini} & \texttt{sim\_shock}, \texttt{sim\_shock\_mese} & --- \\
\bottomrule
\end{tabular}
\caption{Dal simbolo alla cella e alla funzione.}
\end{table}

\section{Il caso di riferimento, svolto}

I valori seguenti sono quelli precaricati nel workbook e verificati con il ricalcolo di Excel. Servono da riscontro: chi rifà i conti a mano su questo caso deve ritrovarli.

\begin{table}[H]
\centering
\small
\begin{tabular}{@{}lrl@{}}
\toprule
Grandezza & Valore & Nota \\
\midrule
Prezzo $P$ & $120\,000$ \euro & acquisto da privato, prima casa, prezzo-valore \\
Rendita catastale $R$ & $450$ \euro & categoria A/3, $55$ m\textsuperscript{2} \\
Valore catastale $V_c$ & $51\,975$ \euro & $450 \cdot 1{,}05 \cdot 110$ \\
Base registro $B$ & $51\,975$ \euro & prezzo-valore attivo \\
Imposte $T_a$ & $1\,140$ \euro & $\max(51\,975\cdot 0{,}02;\,1\,000)+100$ \\
Valore del bonus prima casa & $4\,064$ \euro & differenza col regime ordinario \\
Costi accessori $A$ & $11\,557$ \euro & \\
Costo totale $C$ & $131\,557$ \euro & \\
Incidenza dei costi & $9{,}63$ \% & $A/P$ \\
Esborso $E$ & $41\,557$ \euro & con mutuo di $90\,000$ \\
Rata mensile $\Gamma$ & $436{,}21$ \euro & $i=3{,}2$ \%, $n=300$ \\
Interessi totali & $40\,863$ \euro & \\
Oneri del mutuo $\Omega_M$ & $2\,025$ \euro & \\
Canone lordo $L$ & $6\,000$ \euro & $500$ \euro/mese \\
Ricavo effettivo $L_e$ & $5\,335$ \euro & $\sigma=1$, $\beta=3$ \% \\
$\mathrm{NOI}$ & $1\,805$ \euro & \\
Utile netto $U$ & $684$ \euro & cedolare secca al $21$ \% \\
Rendimento lordo $r_g$ & $5{,}00$ \% & \\
Rendimento netto $r_n$ & $0{,}52$ \% & \\
Cap rate $r_c$ & $1{,}37$ \% & \\
Cash on cash $r_e$ & $-10{,}95$ \% & la leva lavora contro: $r_c < c_d$ \\
$\mathrm{DSCR}$ & $0{,}34$ & sotto uno \\
$\TIR$ & $0{,}40$ \% & orizzonte $25$ anni \\
$\TIR$ reale & $-1{,}56$ \% & con $\pi=2$ \% \\
Rendimento netto reale $r_n^*$ & $-1{,}45$ \% & \\
Prezzo massimo $P^\star$ & $43\,445$ \euro & a $r^\dagger=4$ \% \\
Controllo $\varepsilon$ & $0{,}0000$ \% & \\
\bottomrule
\end{tabular}
\caption{Caso di riferimento con i valori verificati.}
\end{table}

\begin{osservazione}
La lettura del caso è istruttiva e vale enunciarla: con questi numeri l'operazione ha un rendimento nominale positivo ma prossimo a zero, un rendimento reale negativo, un \emph{cash on cash} negativo di undici punti e un DSCR di un terzo. Non è un caso costruito per fallire: è un immobile a $2\,182$ euro il metro quadro affittato a $500$ euro al mese, cioè una configurazione ordinaria in molte zone italiane. La conclusione che il modello suggerisce è che al prezzo trattato l'operazione non regge ai criteri dichiarati, e che il prezzo che li reggerebbe è meno della metà: è esattamente il tipo di risposta per cui uno strumento del genere serve.
\end{osservazione}

\section{Limiti dichiarati}

Si raccolgono qui, in un elenco unico, i limiti che i capitoli hanno introdotto uno per uno.

Il TIR non pesa il rischio e assume il reinvestimento dei flussi allo stesso tasso. La \eqref{eq:leva} vale con costo del debito costante, che in un ammortamento non è. La \eqref{eq:pmax} è esatta solo sul tratto affine del costo, e il minimo di legge dell'imposta di registro lo rompe. Il tornado è univariato e non vede le interazioni. La simulazione estrae variabili indipendenti, mentre nella realtà si muovono insieme. Il piano del mutuo è tabulato su quarant'anni e può non chiudersi entro la tabella. Il confronto comprare-affittare assume disciplina perfetta di chi investe la differenza. La \eqref{eq:scarto} collassa un intervallo sulla sua media. Il modello non prezza il lavoro di gestione oltre al costo figurativo dichiarato, e non modella la ristrutturazione come progetto, per scelta di perimetro.

Nessuno di questi limiti è nascosto in un commento del codice: ciascuno compare nella cella del foglio che lo produce, perché un limite dichiarato dove il numero si legge è un presidio, e un limite dichiarato in una nota tecnica che nessuno apre è un alibi.

\vfill
\noindent\rule{\linewidth}{0.4pt}\\[0.3em]
\small
Questo documento accompagna uno strumento di analisi personale e non costituisce consulenza fiscale, legale o finanziaria. Le aliquote citate sono quelle vigenti alla data di revisione dichiarata nel progetto e cambiano con ogni legge di bilancio. Prima di qualunque firma le posizioni soggettive vanno confermate da un notaio e da un commercialista, e la conformità urbanistica da un tecnico abilitato.

\end{document}
